\section{估值哲学与门禁}
本章发布 AStock 自有综合目标价，不搬运券商目标价，并遵循固定的估值门禁：先分类、再正常化分母、再做情景估值、再构建多锚目标价、再翻译成动作、最后做可复现性审计。核心原则有四条。第一，\textbf{方法匹配业务模型}，绝不把异质产业链套进同一个 PE 表。第二，\textbf{分母必须正常化}，2026E 归母不用单季度机械年化，而用季节性校准并以 TTM（滚动四季度）作防高估地板。第三，\textbf{成长信用必须可追溯}，AI/高速成长溢价要拆到成长段收入与现价隐含增速，近零/负 EPS 标的不给成长 PE 信用。第四，\textbf{市场不是真理也不是噪音}，内在价值锚回答"财务分母支持多少钱"，市场情绪锚回答"市场已经愿意给多少钱以及为什么"，二者按业务模型加权形成综合目标价。所有覆盖标的均具备现价、股本、市值、2026E/2027E/2028E EPS、方法、熊/基准/牛三档、内在锚、市场情绪锚、券商锚、综合目标价、合理区间、隐含空间、动作、催化、失效条件和证据质量。

\begin{exhibitbox}[表：业务模型—估值方法匹配矩阵]
\centering\scriptsize
\begin{tabularx}{\exhibitboxwidth}{>{\bfseries\raggedright\arraybackslash}p{2.6cm} >{\raggedright\arraybackslash\sloppy\hspace{0pt}}p{3.4cm} >{\raggedright\arraybackslash\sloppy\hspace{0pt}}X}
\toprule
\textbf{业务类型} & \textbf{首选方法} & \textbf{必需二级校验} \\
\midrule
AI 光模块龙头 & PE/PEG + 成长盈利桥 & 现价隐含增速、Street 目标价离散度 \\
光芯片/精密器件/稀缺组件 & PE + PS 或稀缺/SOTP 校验 & 客户认证、ASP/毛利率证据 \\
光纤光缆/预制棒（资产型） & 周期正常化 PE + PB/ROE + PS & 经营现金流、运营商/项目订单、光纤价格 \\
网络设备/混合集成 & PE/SOTP 或 PE/PB/PS blend & 在手订单、分部毛利率、现金转化 \\
高速线缆/连接器互连 & PE/PB/PS blend & AI 互连纯度、产品结构、营运资本 \\
近零 EPS / 亏损 & PS、EV/Sales、PB 或观察池 & 盈利路径、稀释/负债敏感性 \\
\bottomrule
\end{tabularx}
\end{exhibitbox}

\section{分母正常化、季节性校准与业绩预告}
把单季度利润乘以 4（或除以固定季节性假设）对周期性和拐点型标的会系统性失真。本报告对每个标的按\textbf{置信度顺序}确定 2026E 分母：\textbf{(1) 管理层 H1 业绩预告（最高置信度前瞻信号）}——已发预告的覆盖标的用 H1 中值除以 H1 占全年比例（0.50）年化得 2026E 归母；因预告理由为\textbf{量价齐升}（量、价同增），2026E 收入也按该利润在 2026Q1 净利率下同步上抬（保守下限，因兑现毛利率实际扩张），使 EPS、PE、PB、PS \textbf{各条估值腿都对预告作出反应}，而不是只动 EPS；(2) 季节性校准 Q1（观测 2025 年 Q1 占全年比例在 8\%--45\% 时用观测值）；(3) TTM（滚动四季度）作防高估地板。已发 H1 预告的名字不能再用陈旧 TTM 或 Q1 年化定分母——例如永鼎股份 2026Q1 归母 1.59 亿元、TTM 仅约 1.03 亿元，但其 H1 预告已达 5.0--7.0 亿元，故 2026E 归母采用预告年化约 12.0 亿元、收入约 94 亿元，EPS 从 0.45 上修到约 0.82，内在价值锚从 14.77 抬升到约 24.57，当前价较内在锚溢价从约 +295\% 收窄到约 +167\%。

\begin{exhibitbox}[表：季节性校准与分母正常化]
\centering\tiny
\renewcommand{\arraystretch}{1.14}
\setlength{\tabcolsep}{2pt}
\begin{tabularx}{\exhibitboxwidth}{>{\bfseries\raggedright\arraybackslash}p{1.4cm} >{\raggedleft\arraybackslash}p{0.82cm} >{\raggedleft\arraybackslash}p{0.82cm} >{\raggedleft\arraybackslash}p{0.82cm} >{\raggedleft\arraybackslash}p{0.82cm} >{\raggedleft\arraybackslash}p{0.9cm} >{\raggedright\arraybackslash\sloppy\hspace{0pt}}X}
\toprule
\textbf{标的} & \textbf{Q1归母} & \textbf{Q1年化} & \textbf{TTM} & \textbf{H1预告} & \textbf{采用口径} & \textbf{说明} \\
\midrule
中际旭创 300308 & 57.35 & 249.33 & 149.49 & -- & 校准 249.33 & 取季节性路径与 TTM 较低者作防高估地板 \\
新易盛 300502 & 27.80 & 120.88 & 107.40 & -- & 校准 120.88 & 取季节性路径与 TTM 较低者作防高估地板 \\
天孚通信 300394 & 4.92 & 20.51 & 21.72 & -- & 校准 20.51 & 取季节性路径与 TTM 较低者作防高估地板 \\
光迅科技 002281 & 2.40 & 10.91 & 10.36 & -- & 校准 10.91 & 取季节性路径与 TTM 较低者作防高估地板 \\
华工科技 000988 & 6.38 & 25.54 & 16.99 & -- & 校准 25.54 & 取季节性路径与 TTM 较低者作防高估地板 \\
剑桥科技 603083 & 1.18 & 5.92 & 3.50 & -- & 校准 5.92 & 取季节性路径与 TTM 较低者作防高估地板 \\
太辰光 300570 & 0.66 & 2.63 & 2.85 & -- & 校准 2.63 & 取季节性路径与 TTM 较低者作防高估地板 \\
源杰科技 688498 & 1.79 & 7.80 & 3.56 & -- & 校准 7.80 & 取季节性路径与 TTM 较低者作防高估地板 \\
亨通光电 600487 & 11.05 & 46.06 & 32.29 & -- & 校准 46.06 & 取季节性路径与 TTM 较低者作防高估地板 \\
中天科技 600522 & 9.19 & 38.29 & 31.94 & -- & 校准 38.29 & 取季节性路径与 TTM 较低者作防高估地板 \\
长飞光纤 601869 & 4.95 & 21.53 & 11.57 & -- & 校准 21.53 & 取季节性路径与 TTM 较低者作防高估地板 \\
烽火通信 600498 & 0.38 & 1.75 & 4.19 & -- & 校准 1.75 & 取季节性路径与 TTM 较低者作防高估地板 \\
中兴通讯 000063 & 13.10 & 59.57 & 44.75 & -- & 校准 59.57 & 取季节性路径与 TTM 较低者作防高估地板 \\
仕佳光子 688313 & 1.16 & 5.05 & 3.95 & -- & 校准 5.05 & 取季节性路径与 TTM 较低者作防高估地板 \\
长光华芯 688048 & 0.04 & 0.22 & 0.34 & -- & 校准 0.22 & 近零 EPS，取季节性与 TTM 较低者，转 PS/PB 校验 \\
光库科技 300620 & 0.45 & 1.95 & 2.11 & -- & 校准 1.95 & 取季节性路径与 TTM 较低者作防高估地板 \\
通鼎互联 002491 & 0.51 & 2.23 & -1.57 & -- & 校准 2.23 & 取季节性路径与 TTM 较低者作防高估地板 \\
永鼎股份 600105 & 1.59 & 6.62 & 1.03 & 12.00 & 预告 12.00 & 采用 H1 业绩预告年化（H1中值/0.50），最高置信度前瞻信号，覆盖 TTM/季节性 \\
长芯博创 300548 & 1.30 & 5.42 & 3.75 & -- & 校准 5.42 & 取季节性路径与 TTM 较低者作防高估地板 \\
德科立 688205 & 0.20 & 0.89 & 0.77 & -- & 校准 0.89 & 取季节性路径与 TTM 较低者作防高估地板 \\
腾景科技 688195 & 0.14 & 0.63 & 0.72 & -- & 校准 0.63 & 取季节性路径与 TTM 较低者作防高估地板 \\
联特科技 301205 & 0.03 & 0.15 & 0.87 & -- & 校准 0.15 & 近零 EPS，取季节性与 TTM 较低者，转 PS/PB 校验 \\
兆龙互连 300913 & 0.49 & 2.02 & 2.47 & -- & 校准 2.02 & 取季节性路径与 TTM 较低者作防高估地板 \\
意华股份 002897 & 0.40 & 1.83 & 2.94 & -- & 校准 1.83 & 取季节性路径与 TTM 较低者作防高估地板 \\
神宇股份 300563 & 0.13 & 0.61 & 0.75 & -- & 校准 0.61 & 取季节性路径与 TTM 较低者作防高估地板 \\
杭电股份 603618 & 0.81 & 3.37 & -2.39 & 7.60 & 预告 7.60 & 采用 H1 业绩预告年化（H1中值/0.50），最高置信度前瞻信号，覆盖 TTM/季节性 \\
\bottomrule
\end{tabularx}
\end{exhibitbox}
\sourcenote{季度-年度桥与三口径：analysis/segment\_forecast\_bridge.md；H1 预告证据：data/earnings\_preview\_h1\_2026\_20260706.md、第 11 章。数据层 data/growth\_driver\_model.json。}

\begin{sourcequalitybox}[业绩预告已计入估值]
本报告 26 只覆盖标的中，有两只已发布 2026 H1 业绩预告，且已按置信度优先原则\textbf{计入其 2026E 归母、2026E 收入、EPS 及 PE/PB/PS 各条估值腿}，而非仅列附录。\textbf{永鼎股份（600105）}：H1 归母 5.0--7.0 亿元、同比 +57\%--+120\%（光纤量价齐升），EPS 0.45$\to$0.82、收入/股 3.55$\to$6.44、内在价值锚 14.77$\to$24.57、溢价 +295\%$\to$+167\%。\textbf{杭电股份（603618）}：H1 归母 3.6--4.0 亿元、同比 +852\%--+958\%（光纤量价齐升，扭亏为盈；FY2025 为亏损），因 Q1 年化与 TTM 都无法代表其扭亏后的盈利能力，采用 H1 预告年化作为 2026E 分母，具体 EPS/内在锚/目标价见其单票估值卡。两只标的综合目标价与评级由各自内在锚、市场情绪锚加权决定：预告显著上修了盈利质量与内在价值锚、收窄了估值泡沫度，但当前价是否给出正空间取决于溢价水平，评级不因单纯利润超预期机械上修。观察池标的锐捷网络（301165）同样已发预告，但不在 26 只当前价目标价覆盖内。第 11 章给出冻结（Q1 年化）与采用（H1 预告）口径的并列对照。价格与未发预告标的的财务分母仍冻结在 2026-06-26。
\end{sourcequalitybox}

\section{最终估值总表}
下表是全部覆盖标的的 AStock 综合目标价、隐含空间与动作。内在价值锚由业务模型组件（PE/PB/PS 按权重）加权得到，综合目标价再叠加市场情绪锚和券商锚。极端"泡沫度"（现价远高于内在锚）不是模型错误，而是当前光通信板块高拥挤定价的真实刻画，报告用市场情绪锚显式承认这部分溢价，同时坚持财务分母底线。

\begin{exhibitbox}[表：AStock 最终估值总表]
\centering
\tiny
\renewcommand{\arraystretch}{1.18}
\setlength{\tabcolsep}{2pt}
\begin{tabularx}{\exhibitboxwidth}{>{\bfseries\raggedright\arraybackslash}p{1.35cm} >{\raggedleft\arraybackslash}p{0.68cm} >{\raggedleft\arraybackslash}p{0.74cm} >{\raggedleft\arraybackslash}p{0.74cm} >{\raggedleft\arraybackslash}p{0.78cm} >{\raggedleft\arraybackslash}p{0.80cm} >{\raggedleft\arraybackslash}p{0.84cm} >{\centering\arraybackslash}p{0.78cm} >{\centering\arraybackslash}p{0.92cm} >{\raggedright\arraybackslash\sloppy\hspace{0pt}}X}
\toprule
\textbf{标的} & \textbf{现价} & \textbf{市值} & \textbf{26E EPS} & \textbf{27E EPS} & \textbf{综合目标} & \textbf{空间} & \textbf{动作} & \textbf{方法/证据} & \textbf{估值解释} \\
\midrule
\rowcolor{navy!6}中际旭创 300308 & 1253.89 & 13881 & 22.52 & 29.73 & 1229.56 & -1.9\% & \stance{riskamber}{中性} & \makecell[c]{PE/PEG\\A-} & 800G/1.6T 规模和盈利兑现最强，但当前估值已经反映海外 AI 需求持续高景气。 \\
\rowcolor{navy!6}新易盛 300502 & 566.00 & 5620 & 12.17 & 15.83 & 576.55 & +1.9\% & \stance{riskamber}{中性} & \makecell[c]{PE/PEG\\A-} & 高增长光模块弹性强，2025--2026 年交付较强，但估值对单一客户和 ASP 假设敏感。 \\
\rowcolor{accentblue!6}天孚通信 300394 & 318.00 & 2472 & 2.64 & 3.27 & 254.40 & -20.0\% & \stance{riskamber}{中性观察} & \makecell[c]{PE/PEG\\B+} & 高毛利精密器件平台价值明确，但增速相对模块龙头放缓。 \\
\rowcolor{accentblue!6}光迅科技 002281 & 240.34 & 1922 & 1.36 & 1.64 & 187.47 & -22.0\% & \stance{riskamber}{中性观察} & \makecell[c]{PE/PB/PS\\B} & 产品线和央企平台优势存在，但盈利能力显著低于 AI 光模块龙头。 \\
\rowcolor{riskamber!8}华工科技 000988 & 160.94 & 1606 & 2.56 & 3.02 & 128.75 & -20.0\% & \stance{riskamber}{中性观察} & \makecell[c]{PE/PB/PS\\B} & 混合业务提供盈利稳定性，但削弱纯 AI 光模块弹性。 \\
\rowcolor{riskamber!8}剑桥科技 603083 & 245.92 & 856 & 1.70 & 2.21 & 177.06 & -28.0\% & \stance{riskamber}{中性观察} & \makecell[c]{PE/PB/PS\\C+} & 反转弹性存在，但 Q1 利润分母仍偏小，难以支撑过高目标价。 \\
\rowcolor{riskamber!8}太辰光 300570 & 234.10 & 531 & 1.16 & 1.33 & 163.87 & -30.0\% & \stance{riskamber}{中性观察} & \makecell[c]{PE/PB/PS\\C} & 无源器件敞口真实，但 2026Q1 收入/利润下滑，估值需要交付修复。 \\
\rowcolor{riskamber!8}源杰科技 688498 & 1776.91 & 1526 & 9.09 & 12.54 & 1421.53 & -20.0\% & \stance{riskamber}{中性观察} & \makecell[c]{PE/PB/PS\\B} & 上游战略属性最强之一，但当前价格已经前置长期兑现路径。 \\
\rowcolor{accentblue!6}亨通光电 600487 & 110.81 & 2684 & 1.90 & 2.13 & 86.43 & -22.0\% & \stance{riskamber}{中性观察} & \makecell[c]{PE/PB/PS\\B} & 光纤光缆龙头兼具海缆和电网敞口，但 AI 数据中心纯度低于模块龙头。 \\
\rowcolor{accentblue!6}中天科技 600522 & 61.29 & 2071 & 1.13 & 1.27 & 47.81 & -22.0\% & \stance{riskamber}{中性观察} & \makecell[c]{PE/PB/PS\\B} & 电力海洋与通信线缆平台提供盈利韧性，但不能按纯 AI 光模块估值。 \\
\rowcolor{accentblue!6}长飞光纤 601869 & 543.00 & 4481 & 2.61 & 3.03 & 434.40 & -20.0\% & \stance{riskamber}{中性观察} & \makecell[c]{PE/PB/PS\\B+} & 预制棒/光纤/光缆一体化稀缺性较强，但数据中心需求需体现在利用率上。 \\
\rowcolor{accentblue!6}烽火通信 600498 & 76.22 & 975 & 0.14 & 0.16 & 53.35 & -30.0\% & \stance{riskamber}{中性观察} & \makecell[c]{PE/PB/PS\\B-} & 运营商光网络平台具备战略意义，但 Q1 利润分母薄，周期慢于云侧 AI capex。 \\
\rowcolor{accentblue!6}中兴通讯 000063 & 35.67 & 1731 & 1.23 & 1.37 & 31.06 & -12.9\% & \stance{riskamber}{中性} & \makecell[c]{PE/PB/PS\\A-} & 覆盖运营商、云网络和应用层，但光通信只是业务组合的一部分。 \\
\rowcolor{riskamber!8}仕佳光子 688313 & 179.70 & 812 & 1.12 & 1.40 & 125.79 & -30.0\% & \stance{riskamber}{中性观察} & \makecell[c]{PE/PB/PS\\B} & 无源光芯片与 WDM/AWG/硅光生态相关性强，但估值需要高利用率支撑。 \\
\rowcolor{riskamber!8}长光华芯 688048 & 396.00 & 698 & 0.13 & 0.18 & 277.20 & -30.0\% & \stance{riskamber}{中性观察} & \makecell[c]{PE/PB/PS\\C+} & 激光芯片期权高，但 Q1 利润分母仍小，基准情景要求快速认证和放量。 \\
\rowcolor{accentblue!6}光库科技 300620 & 362.43 & 903 & 0.78 & 0.98 & 253.70 & -30.0\% & \stance{riskamber}{中性观察} & \makecell[c]{PE/PB/PS\\B} & 光器件和薄膜铌酸锂期权提高战略相关性，但股价已经计入较长兑现路径。 \\
\rowcolor{accentblue!6}通鼎互联 002491 & 36.15 & 444 & 0.18 & 0.20 & 26.03 & -28.0\% & \stance{riskamber}{中性观察} & \makecell[c]{PE/PB/PS\\B-} & 属于完整产业链底座，但盈利质量和 AI 数据中心纯度低于模块龙头。 \\
\rowcolor{accentblue!6}永鼎股份 600105 & 65.71 & 961 & 0.82 & 0.90 & 46.00 & -30.0\% & \stance{riskamber}{中性观察} & \makecell[c]{PE/PB/PS\\B} & 线缆和电信基础设施修复可见，但业务慢周期、纯度低于 AI 光模块。 \\
\rowcolor{accentblue!6}长芯博创 300548 & 262.65 & 759 & 1.88 & 2.34 & 189.11 & -28.0\% & \stance{riskamber}{中性观察} & \makecell[c]{PE/PB/PS\\B} & 光模块/器件平台具备数通敞口，但当前估值要求高速产品持续交付。 \\
\rowcolor{riskamber!8}德科立 688205 & 210.52 & 345 & 0.55 & 0.68 & 73.50 & -65.1\% & \stance{riskred}{减持} & \makecell[c]{PE/PB/PS\\B} & 补足 DCI 和运营商相干传输链条，但 Q1 利润基数仍偏小。 \\
\rowcolor{accentblue!6}腾景科技 688195 & 187.60 & 246 & 0.48 & 0.57 & 83.18 & -55.7\% & \stance{riskred}{减持} & \makecell[c]{PE/PB/PS\\B} & 精密光学与高速光系统相关，但规模小，估值需要持续订单转化。 \\
\rowcolor{riskamber!8}联特科技 301205 & 315.86 & 409 & 0.12 & 0.16 & 116.17 & -63.2\% & \stance{riskred}{减持} & \makecell[c]{PE/PB/PS\\C+} & 高速模块弹性大，但 Q1 利润分母很小，目标价高度依赖执行。 \\
\rowcolor{riskamber!8}兆龙互连 300913 & 50.75 & 174 & 0.59 & 0.70 & 22.44 & -55.8\% & \stance{riskred}{减持} & \makecell[c]{PE/PB/PS\\B} & 高速线缆/铜互连补足机柜内和企业网络层，但不是纯光模块替代品。 \\
\rowcolor{riskamber!8}意华股份 002897 & 84.50 & 162 & 0.95 & 1.10 & 49.92 & -40.9\% & \stance{riskred}{减持} & \makecell[c]{PE/PB/PS\\B-} & 连接器敞口属于数据中心互连链，但混合业务需要纯度折价。 \\
\rowcolor{riskamber!8}神宇股份 300563 & 24.73 & 48 & 0.32 & 0.37 & 12.95 & -47.7\% & \stance{riskred}{减持} & \makecell[c]{PE/PB/PS\\C+} & 线缆/互连拓宽产业链覆盖，但光通信纯度和 Q1 利润规模有限。 \\
\rowcolor{accentblue!6}杭电股份 603618 & 52.00 & 350 & 1.13 & 1.30 & 25.80 & -50.4\% & \stance{riskred}{减持} & \makecell[c]{PE/PB/PS\\C+} & AI 驱动光纤需求带动 H1 2026 从 2025 年亏损大幅扭亏，但基础业务为低毛利线缆，复苏需在一次预告之外持续确认。 \\
\bottomrule
\end{tabularx}
\end{exhibitbox}
\sourcenote{估值模型：data/current\_valuation\_model\_20260626.json；可复现性见 analysis/valuation\_audit.md（Model Reproducibility: PASS）。}

\section{成长盈利拆分与现价隐含增速}
AI/高速成长溢价必须回答"贵在哪一段、现价隐含多快增速"。本节把 2026E 收入拆成基础业务段与 AI/高速成长段（AI 收入占比为基于业务层级的 AStock 建模假设，非公司分部披露），并用反向估值给出：在 AStock 基准 PE 档下，当前价要求未来 3 年归母 CAGR 达到多少，再与 AStock 假设的 2026--2028 增速对比。隐含增速远高于假设=估值把成长前置；接近=需持续兑现；低于=若兑现有修复空间。近零/负 EPS 标的（如长光华芯）无法用 PE 反推，转 PS/PB 并降为观察池，不给成长 PE 信用。

\begin{exhibitbox}[表：成长盈利拆分与现价隐含增速]
\centering\tiny
\renewcommand{\arraystretch}{1.14}
\setlength{\tabcolsep}{2pt}
\begin{tabularx}{\exhibitboxwidth}{>{\bfseries\raggedright\arraybackslash}p{1.5cm} >{\raggedleft\arraybackslash}p{0.95cm} >{\raggedleft\arraybackslash}p{1.05cm} >{\raggedleft\arraybackslash}p{1.05cm} >{\raggedleft\arraybackslash}p{1.05cm} >{\raggedleft\arraybackslash}p{1.05cm} >{\raggedright\arraybackslash\sloppy\hspace{0pt}}X}
\toprule
\textbf{标的} & \textbf{AI占比} & \textbf{成长段归母} & \textbf{假设增速} & \textbf{隐含CAGR} & \textbf{隐含-假设} & \textbf{成长估值信用与判读} \\
\midrule
中际旭创 300308 & 90\% & 224.39 & 27\% & 7\% & -20pct & 成长 PE/PEG；若兑现有修复空间：隐含低于假设 \\
新易盛 300502 & 90\% & 108.79 & 26\% & 5\% & -21pct & 成长 PE/PEG；若兑现有修复空间：隐含低于假设 \\
天孚通信 300394 & 72\% & 14.77 & 21\% & 39\% & 18pct & 成长 PE/PEG；估值前置：现价隐含增速高于假设 \\
光迅科技 002281 & 70\% & 7.63 & 17\% & 71\% & 54pct & 成长 PE/PEG；估值前置：现价隐含增速高于假设 \\
华工科技 000988 & 45\% & 11.49 & 16\% & 22\% & 6pct & 混合有限成长；估值前置：现价隐含增速高于假设 \\
剑桥科技 603083 & 70\% & 4.14 & 27\% & 48\% & 20pct & 成长 PE/PEG；估值前置：现价隐含增速高于假设 \\
太辰光 300570 & 55\% & 1.45 & 13\% & 79\% & 66pct & 混合有限成长；估值前置：现价隐含增速高于假设 \\
源杰科技 688498 & 65\% & 5.07 & 34\% & 44\% & 10pct & 成长 PE/PEG；估值前置：现价隐含增速高于假设 \\
亨通光电 600487 & 25\% & 11.51 & 11\% & 34\% & 23pct & 混合有限成长；估值前置：现价隐含增速高于假设 \\
中天科技 600522 & 25\% & 9.57 & 11\% & 31\% & 20pct & 混合有限成长；估值前置：现价隐含增速高于假设 \\
长飞光纤 601869 & 25\% & 5.38 & 14\% & 95\% & 81pct & 混合有限成长；估值前置：现价隐含增速高于假设 \\
烽火通信 600498 & 30\% & 0.52 & 13\% & 171\% & 158pct & 混合有限成长；估值前置：现价隐含增速高于假设 \\
中兴通讯 000063 & 30\% & 17.87 & 11\% & 13\% & 2pct & 混合有限成长；需持续兑现：隐含≈假设 \\
仕佳光子 688313 & 65\% & 3.28 & 21\% & 48\% & 26pct & 成长 PE/PEG；估值前置：现价隐含增速高于假设 \\
长光华芯 688048 & 65\% & 0.15 & 40\% & 254\% & 215pct & 观察池；估值前置：现价隐含增速高于假设 \\
光库科技 300620 & 60\% & 1.17 & 21\% & 110\% & 89pct & 成长 PE/PEG；估值前置：现价隐含增速高于假设 \\
通鼎互联 002491 & 25\% & 0.56 & 11\% & 103\% & 92pct & 混合有限成长；估值前置：现价隐含增速高于假设 \\
永鼎股份 600105 & 25\% & 3.00 & 9\% & 52\% & 43pct & 混合有限成长；估值前置：现价隐含增速高于假设 \\
长芯博创 300548 & 70\% & 3.79 & 21\% & 49\% & 28pct & 成长 PE/PEG；估值前置：现价隐含增速高于假设 \\
德科立 688205 & 70\% & 0.62 & 21\% & 105\% & 84pct & 成长 PE/PEG；估值前置：现价隐含增速高于假设 \\
腾景科技 688195 & 72\% & 0.45 & 17\% & 111\% & 93pct & 成长 PE/PEG；估值前置：现价隐含增速高于假设 \\
联特科技 301205 & 70\% & 0.11 & 30\% & 266\% & 236pct & 观察池；估值前置：现价隐含增速高于假设 \\
兆龙互连 300913 & 45\% & 0.91 & 16\% & 35\% & 19pct & 混合有限成长；估值前置：现价隐含增速高于假设 \\
意华股份 002897 & 45\% & 0.82 & 13\% & 43\% & 30pct & 混合有限成长；估值前置：现价隐含增速高于假设 \\
神宇股份 300563 & 45\% & 0.28 & 13\% & 34\% & 21pct & 混合有限成长；估值前置：现价隐含增速高于假设 \\
杭电股份 603618 & 25\% & 1.90 & 12\% & 28\% & 15pct & 混合有限成长；估值前置：现价隐含增速高于假设 \\
\bottomrule
\end{tabularx}
\end{exhibitbox}
\sourcenote{成长盈利模型：analysis/growth\_earnings\_model.md、analysis/implied\_growth\_sensitivity.md、data/growth\_driver\_model.json。AI 收入占比与增速为 AStock 建模假设，非公司财报分部披露。}

\section{基于 2026E 收入和成长性的市场预期估值}
本节补上"市场炒预期"的视角。AStock 内在价值锚回答按业务模型和 2027E 正常化分母的冷静价值；市场预期估值回答：如果投资者愿意基于 2026E 收入、2026E EPS 和 2027E 成长性给更高预期倍数，当前价格对应的上方或下方空间是多少。预期价值不是买入理由，只有当下一季度收入、毛利率、现金流和订单质量继续验证，预期倍数才有资格维持。

\begin{exhibitbox}[表：市场预期估值桥]
\centering\tiny
\renewcommand{\arraystretch}{1.14}
\setlength{\tabcolsep}{2pt}
\begin{tabularx}{\exhibitboxwidth}{>{\bfseries\raggedright\arraybackslash}p{1.25cm} >{\raggedleft\arraybackslash}p{0.72cm} >{\raggedleft\arraybackslash}p{0.82cm} >{\raggedleft\arraybackslash}p{0.82cm} >{\raggedleft\arraybackslash}p{0.78cm} >{\raggedright\arraybackslash\sloppy\hspace{0pt}}p{1.65cm} >{\raggedleft\arraybackslash}p{0.82cm} >{\raggedleft\arraybackslash}p{0.82cm} >{\raggedright\arraybackslash\sloppy\hspace{0pt}}X}
\toprule
\textbf{标的} & \textbf{现价} & \textbf{2026E收入} & \textbf{收入增} & \textbf{2026E EPS} & \textbf{预期倍数} & \textbf{预期价} & \textbf{空间} & \textbf{预期驱动} \\
\midrule
中际旭创 300308 & 1253.89 & 847.67 & +32.0\% & 22.52 & PE 51.5x / PS 0.0x & 1159.42 & -7.5\% & 收入增长和毛利率兑现推动 EPS 预期，核心风险是海外云厂商订单持续性。 \\
新易盛 300502 & 566.00 & 362.52 & +30.0\% & 12.17 & PE 45.4x / PS 0.0x & 552.70 & -2.4\% & 收入增长和毛利率兑现推动 EPS 预期，核心风险是海外云厂商订单持续性。 \\
天孚通信 300394 & 318.00 & 55.43 & +24.0\% & 2.64 & PE 49.9x / PS 0.0x & 131.53 & -58.6\% & 收入增长和毛利率兑现推动 EPS 预期，核心风险是海外云厂商订单持续性。 \\
光迅科技 002281 & 240.34 & 126.06 & +20.0\% & 1.36 & PE 38.2x / PS 4.0x & 54.27 & -77.4\% & 市场预期来自收入增长、毛利率改善和业务纯度提升。 \\
华工科技 000988 & 160.94 & 170.63 & +18.0\% & 2.56 & PE 37.8x / PS 3.9x & 86.03 & -46.5\% & 市场预期来自收入增长、毛利率改善和业务纯度提升。 \\
剑桥科技 603083 & 245.92 & 64.33 & +30.0\% & 1.70 & PE 51.1x / PS 6.0x & 97.08 & -60.5\% & 市场预期主要来自高增长和战略稀缺性，需用客户认证和收入规模验证。 \\
太辰光 300570 & 234.10 & 13.63 & +15.0\% & 1.16 & PE 37.4x / PS 0.9x & 25.19 & -89.2\% & 收入增长预期主要来自运营商/云侧线缆项目，估值上限受现金流和线缆价格约束。 \\
源杰科技 688498 & 1776.91 & 15.45 & +38.0\% & 9.09 & PE 76.1x / PS 6.2x & 375.41 & -78.9\% & 市场预期主要来自高增长和战略稀缺性，需用客户认证和收入规模验证。 \\
亨通光电 600487 & 110.81 & 741.27 & +12.0\% & 1.90 & PE 25.3x / PS 0.9x & 35.55 & -67.9\% & 收入增长预期主要来自运营商/云侧线缆项目，估值上限受现金流和线缆价格约束。 \\
中天科技 600522 & 61.29 & 547.59 & +12.0\% & 1.13 & PE 25.3x / PS 0.9x & 22.59 & -63.1\% & 收入增长预期主要来自运营商/云侧线缆项目，估值上限受现金流和线缆价格约束。 \\
长飞光纤 601869 & 543.00 & 160.67 & +16.0\% & 2.61 & PE 30.0x / PS 0.9x & 49.19 & -90.9\% & 收入增长预期主要来自运营商/云侧线缆项目，估值上限受现金流和线缆价格约束。 \\
烽火通信 600498 & 76.22 & 205.01 & +15.0\% & 0.14 & PE 29.9x / PS 1.1x & 12.27 & -83.9\% & 收入增长预期取决于运营商资本开支、云网络设备和非光通信业务结构。 \\
中兴通讯 000063 & 35.67 & 1590.37 & +12.0\% & 1.23 & PE 21.1x / PS 1.1x & 30.00 & -15.9\% & 收入增长预期取决于运营商资本开支、云网络设备和非光通信业务结构。 \\
仕佳光子 688313 & 179.70 & 25.07 & +25.0\% & 1.12 & PE 55.6x / PS 5.8x & 43.20 & -76.0\% & 市场预期来自收入增长、毛利率改善和业务纯度提升。 \\
长光华芯 688048 & 396.00 & 6.50 & +45.0\% & 0.13 & PE 84.2x / PS 6.5x & 25.74 & -93.5\% & 市场预期主要来自高增长和战略稀缺性，需用客户认证和收入规模验证。 \\
光库科技 300620 & 362.43 & 18.54 & +25.0\% & 0.78 & PE 55.6x / PS 0.9x & 26.34 & -92.7\% & 收入增长预期主要来自运营商/云侧线缆项目，估值上限受现金流和线缆价格约束。 \\
通鼎互联 002491 & 36.15 & 46.24 & +12.0\% & 0.18 & PE 25.3x / PS 5.4x & 13.59 & -62.4\% & 收入增长预期主要来自运营商/云侧线缆项目，估值上限受现金流和线缆价格约束。 \\
永鼎股份 600105 & 65.71 & 94.12 & +10.0\% & 0.82 & PE 24.0x / PS 5.3x & 23.87 & -63.7\% & 收入增长预期主要来自运营商/云侧线缆项目，估值上限受现金流和线缆价格约束。 \\
长芯博创 300548 & 262.65 & 27.94 & +25.0\% & 1.88 & PE 46.7x / PS 4.1x & 72.24 & -72.5\% & 市场预期来自收入增长、毛利率改善和业务纯度提升。 \\
德科立 688205 & 210.52 & 11.57 & +24.0\% & 0.55 & PE 49.9x / PS 5.8x & 37.98 & -82.0\% & 市场预期来自收入增长、毛利率改善和业务纯度提升。 \\
腾景科技 688195 & 187.60 & 7.44 & +20.0\% & 0.48 & PE 45.8x / PS 5.6x & 27.63 & -85.3\% & 市场预期来自收入增长、毛利率改善和业务纯度提升。 \\
联特科技 301205 & 315.86 & 10.67 & +35.0\% & 0.12 & PE 63.7x / PS 6.1x & 31.76 & -89.9\% & 市场预期主要来自高增长和战略稀缺性，需用客户认证和收入规模验证。 \\
兆龙互连 300913 & 50.75 & 21.19 & +18.0\% & 0.59 & PE 37.8x / PS 2.0x & 20.34 & -59.9\% & 收入增长预期来自高速线缆/连接器放量，但需要产品结构和营运资本验证。 \\
意华股份 002897 & 84.50 & 58.75 & +15.0\% & 0.95 & PE 32.0x / PS 2.0x & 46.72 & -44.7\% & 收入增长预期来自高速线缆/连接器放量，但需要产品结构和营运资本验证。 \\
神宇股份 300563 & 24.73 & 9.25 & +15.0\% & 0.32 & PE 34.2x / PS 2.0x & 13.41 & -45.8\% & 收入增长预期来自高速线缆/连接器放量，但需要产品结构和营运资本验证。 \\
杭电股份 603618 & 52.00 & 207.05 & +15.0\% & 1.13 & PE 23.5x / PS 0.9x & 20.86 & -59.9\% & 收入增长预期主要来自运营商/云侧线缆项目，估值上限受现金流和线缆价格约束。 \\
\bottomrule
\end{tabularx}
\end{exhibitbox}
\sourcenote{预期估值模型：data/market\_expectation\_valuation\_20260626.json。}

\section{市场隐含预期与情绪锚}
内在价值锚回答"财务分母支持多少钱"，市场隐含预期锚回答"当前市场已经愿意给多少钱以及为什么"。本节用成交额分位、当前隐含 PE/PS/PB、券商锚和情绪状态修正综合目标价。它不是把现价当真理，而是防止模型在强共识阶段把慢周期或期权型标的机械低估；反过来，当内在锚远低于现价时，报告也不机械写减持，而是标注情绪溢价并给出失效条件。

\begin{exhibitbox}[表：市场隐含预期与情绪锚]
\centering\tiny
\renewcommand{\arraystretch}{1.13}
\setlength{\tabcolsep}{2pt}
\begin{tabularx}{\exhibitboxwidth}{>{\bfseries\raggedright\arraybackslash}p{1.25cm} >{\raggedleft\arraybackslash}p{0.72cm} >{\raggedleft\arraybackslash}p{0.78cm} >{\raggedright\arraybackslash\sloppy\hspace{0pt}}p{1.45cm} >{\raggedleft\arraybackslash}p{0.72cm} >{\centering\arraybackslash}p{0.78cm} >{\raggedleft\arraybackslash}p{0.78cm} >{\raggedleft\arraybackslash}p{0.78cm} >{\raggedright\arraybackslash\sloppy\hspace{0pt}}p{1.25cm} >{\raggedleft\arraybackslash}p{0.82cm} >{\raggedright\arraybackslash\sloppy\hspace{0pt}}X}
\toprule
\textbf{标的} & \textbf{内在锚} & \textbf{市场锚} & \textbf{当前隐含倍数} & \textbf{成交分位} & \textbf{情绪} & \textbf{券商锚} & \textbf{综合目标} & \textbf{权重} & \textbf{空间} & \textbf{动作逻辑} \\
\midrule
中际旭创 300308 & 1337.79 & 1003.11 & PE 55.7倍 / PS 16.4倍 & +100.0\% & 强共识/高拥挤 & 1062.50 & 1229.56 & 内在65\% / 市场20\% / 券商15\% & -1.9\% & 动作由综合目标价相对当前价决定。 \\
新易盛 300502 & 633.04 & 452.80 & PE 46.5倍 / PS 15.5倍 & +96.2\% & 强共识/高拥挤 & 496.76 & 576.55 & 内在65\% / 市场20\% / 券商15\% & +1.9\% & 动作由综合目标价相对当前价决定。 \\
天孚通信 300394 & 147.20 & 254.40 & PE 120.5倍 / PS 44.6倍 & +84.6\% & 强共识/高拥挤 & 286.81 & 254.40 & 内在57\% / 市场28\% / 券商15\% & -20.0\% & 市场高成交已给出情绪溢价，维持中性观察并等待利润/现金流追认。 \\
光迅科技 002281 & 58.46 & 187.47 & PE 176.2倍 / PS 15.2倍 & +80.8\% & 强共识/高拥挤 & 123.86 & 187.47 & 内在47\% / 市场38\% / 券商15\% & -22.0\% & 市场高成交已给出情绪溢价，维持中性观察并等待利润/现金流追认。 \\
华工科技 000988 & 92.79 & 128.75 & PE 62.9倍 / PS 9.4倍 & +76.9\% & 强共识/高拥挤 & 133.99 & 128.75 & 内在47\% / 市场38\% / 券商15\% & -20.0\% & 市场高成交已给出情绪溢价，维持中性观察并等待利润/现金流追认。 \\
剑桥科技 603083 & 105.68 & 177.06 & PE 144.7倍 / PS 13.3倍 & +42.3\% & 活跃共识 & 165.68 & 177.06 & 内在47\% / 市场38\% / 券商15\% & -28.0\% & 市场高成交已给出情绪溢价，维持中性观察并等待利润/现金流追认。 \\
太辰光 300570 & 26.60 & 163.87 & PE 201.8倍 / PS 39.0倍 & +53.8\% & 活跃共识 & 148.32 & 163.87 & 内在38\% / 市场47\% / 券商15\% & -30.0\% & 市场高成交已给出情绪溢价，维持中性观察并等待利润/现金流追认。 \\
源杰科技 688498 & 434.08 & 1421.53 & PE 195.5倍 / PS 98.8倍 & +65.4\% & 强共识/高拥挤 & 1164.00 & 1421.53 & 内在47\% / 市场38\% / 券商15\% & -20.0\% & 市场高成交已给出情绪溢价，维持中性观察并等待利润/现金流追认。 \\
亨通光电 600487 & 36.95 & 86.43 & PE 58.3倍 / PS 3.6倍 & +88.5\% & 强共识/高拥挤 & 97.88 & 86.43 & 内在38\% / 市场47\% / 券商15\% & -22.0\% & 市场高成交已给出情绪溢价，维持中性观察并等待利润/现金流追认。 \\
中天科技 600522 & 23.35 & 47.81 & PE 54.1倍 / PS 3.8倍 & +92.3\% & 强共识/高拥挤 & 49.61 & 47.81 & 内在38\% / 市场47\% / 券商15\% & -22.0\% & 市场高成交已给出情绪溢价，维持中性观察并等待利润/现金流追认。 \\
长飞光纤 601869 & 51.93 & 434.40 & PE 208.2倍 / PS 27.9倍 & +61.5\% & 强共识/高拥挤 & 292.37 & 434.40 & 内在38\% / 市场47\% / 券商15\% & -20.0\% & 市场高成交已给出情绪溢价，维持中性观察并等待利润/现金流追认。 \\
烽火通信 600498 & 12.45 & 53.35 & PE 558.9倍 / PS 4.8倍 & +69.2\% & 活跃共识 & 未披露 & 53.35 & 内在45\% / 市场55\% / 券商0\% & -30.0\% & 市场高成交已给出情绪溢价，维持中性观察并等待利润/现金流追认。 \\
中兴通讯 000063 & 31.03 & 25.68 & PE 29.1倍 / PS 1.1倍 & +57.7\% & 活跃共识 & 43.69 & 31.06 & 内在50\% / 市场35\% / 券商15\% & -12.9\% & 动作由综合目标价相对当前价决定。 \\
仕佳光子 688313 & 47.49 & 125.79 & PE 160.8倍 / PS 32.4倍 & +46.2\% & 活跃共识 & 157.53 & 125.79 & 内在47\% / 市场38\% / 券商15\% & -30.0\% & 市场高成交已给出情绪溢价，维持中性观察并等待利润/现金流追认。 \\
长光华芯 688048 & 26.74 & 277.20 & PE 3118.1倍 / PS 107.5倍 & +34.6\% & 活跃共识 & 86.84 & 277.20 & 内在47\% / 市场38\% / 券商15\% & -30.0\% & 市场高成交已给出情绪溢价，维持中性观察并等待利润/现金流追认。 \\
光库科技 300620 & 28.59 & 253.70 & PE 464.4倍 / PS 48.7倍 & +30.8\% & 活跃共识 & 177.00 & 253.70 & 内在38\% / 市场47\% / 券商15\% & -30.0\% & 市场高成交已给出情绪溢价，维持中性观察并等待利润/现金流追认。 \\
通鼎互联 002491 & 13.96 & 26.03 & PE 199.4倍 / PS 9.6倍 & +50.0\% & 活跃共识 & 未披露 & 26.03 & 内在45\% / 市场55\% / 券商0\% & -28.0\% & 市场高成交已给出情绪溢价，维持中性观察并等待利润/现金流追认。 \\
永鼎股份 600105 & 24.57 & 46.00 & PE 80.1倍 / PS 10.2倍 & +73.1\% & 活跃共识 & 未披露 & 46.00 & 内在45\% / 市场55\% / 券商0\% & -30.0\% & 市场高成交已给出情绪溢价，维持中性观察并等待利润/现金流追认。 \\
长芯博创 300548 & 80.23 & 189.11 & PE 140.1倍 / PS 27.2倍 & +38.5\% & 活跃共识 & 未披露 & 189.11 & 内在55\% / 市场45\% / 券商0\% & -28.0\% & 市场高成交已给出情绪溢价，维持中性观察并等待利润/现金流追认。 \\
德科立 688205 & 40.10 & 134.73 & PE 386.0倍 / PS 29.8倍 & +15.4\% & 中性交易 & 未披露 & 73.50 & 内在65\% / 市场35\% / 券商0\% & -65.1\% & 动作由综合目标价相对当前价决定。 \\
腾景科技 688195 & 29.18 & 112.56 & PE 392.3倍 / PS 33.1倍 & +19.2\% & 中性交易 & 222.42 & 83.18 & 内在55\% / 市场30\% / 券商15\% & -55.7\% & 动作由综合目标价相对当前价决定。 \\
联特科技 301205 & 33.68 & 189.52 & PE 2688.2倍 / PS 38.4倍 & +23.1\% & 中性交易 & 271.98 & 116.17 & 内在55\% / 市场30\% / 券商15\% & -63.2\% & 动作由综合目标价相对当前价决定。 \\
兆龙互连 300913 & 21.10 & 24.36 & PE 86.0倍 / PS 8.2倍 & +7.7\% & 弱共识 & 未披露 & 22.44 & 内在59\% / 市场41\% / 券商0\% & -55.8\% & 动作由综合目标价相对当前价决定。 \\
意华股份 002897 & 48.20 & 52.39 & PE 88.5倍 / PS 2.8倍 & +11.5\% & 中性交易 & 未披露 & 49.92 & 内在59\% / 市场41\% / 券商0\% & -40.9\% & 动作由综合目标价相对当前价决定。 \\
神宇股份 300563 & 13.70 & 11.87 & PE 77.7倍 / PS 5.2倍 & +3.8\% & 弱共识 & 未披露 & 12.95 & 内在59\% / 市场41\% / 券商0\% & -47.7\% & 动作由综合目标价相对当前价决定。 \\
杭电股份 603618 & 22.03 & 31.20 & PE 46.1倍 / PS 1.7倍 & +26.9\% & 中性交易 & 未披露 & 25.80 & 内在59\% / 市场41\% / 券商0\% & -50.4\% & 动作由综合目标价相对当前价决定。 \\
\bottomrule
\end{tabularx}
\end{exhibitbox}
\sourcenote{市场情绪锚：data/market\_sentiment\_anchor\_20260626.json；方法论参考 CFA 市场法估值、预期投资反推框架和投资者情绪研究。}

\section{券商/Street 对照与发布降级}
估值门禁要求全产业链报告为每个覆盖标的提供逐篇原文可复核的券商目标价与 2026E/2027E 预测。本轮多数标的仅取得第三方一致预期页面或媒体摘要，部分标的无可用券商目标价。按门禁，券商原文覆盖不完整时不得给 final\_signoff: PASS，故本报告显式降级为 \textbf{MECHANICAL\_PASS\_INSTITUTIONAL\_FAIL}：AStock 自有目标价、内在锚、市场情绪锚与成长盈利模型可独立复算并支撑全部结论，但券商对照层为机构级不完整，任何投资结论不得单独依赖券商锚。预测收入/净利润/EPS 与估值方法字段在公开页面未逐项披露，统一标注 not disclosed，且券商锚仅在证据质量足够时进入权重。

\begin{exhibitbox}[表：券商/Street 对照与证据质量]
\centering\tiny
\renewcommand{\arraystretch}{1.14}
\setlength{\tabcolsep}{2pt}
\begin{tabularx}{\exhibitboxwidth}{>{\bfseries\raggedright\arraybackslash}p{1.5cm} >{\raggedright\arraybackslash\sloppy\hspace{0pt}}p{2.4cm} >{\raggedleft\arraybackslash}p{0.95cm} >{\raggedleft\arraybackslash}p{0.95cm} >{\raggedleft\arraybackslash}p{0.85cm} >{\centering\arraybackslash}p{1.0cm} >{\raggedright\arraybackslash\sloppy\hspace{0pt}}X}
\toprule
\textbf{标的} & \textbf{来源类型} & \textbf{目标价} & \textbf{隐含空间} & \textbf{券商权重} & \textbf{证据质量} & \textbf{预测明细} \\
\midrule
中际旭创 300308 & 第三方一致预期页面 & 1062.50 & -15.3\% & 15\% & B & 收入/利润/EPS: not disclosed \\
新易盛 300502 & 第三方一致预期页面 & 496.76 & -12.2\% & 15\% & B & 收入/利润/EPS: not disclosed \\
天孚通信 300394 & 第三方一致预期页面 & 286.81 & -9.8\% & 15\% & B & 收入/利润/EPS: not disclosed \\
光迅科技 002281 & 第三方一致预期页面 & 123.86 & -48.5\% & 15\% & B & 收入/利润/EPS: not disclosed \\
华工科技 000988 & 第三方一致预期页面 & 133.99 & -16.7\% & 15\% & B- & 收入/利润/EPS: not disclosed \\
剑桥科技 603083 & 第三方一致预期页面 & 165.68 & -32.6\% & 15\% & C+ & 收入/利润/EPS: not disclosed \\
太辰光 300570 & 第三方一致预期页面 & 148.32 & -36.6\% & 15\% & C+ & 收入/利润/EPS: not disclosed \\
源杰科技 688498 & 第三方一致预期页面 & 1164.00 & -34.5\% & 15\% & B- & 收入/利润/EPS: not disclosed \\
亨通光电 600487 & 第三方一致预期页面 & 97.88 & -11.7\% & 15\% & B & 收入/利润/EPS: not disclosed \\
中天科技 600522 & 第三方一致预期页面 & 49.61 & -19.1\% & 15\% & B & 收入/利润/EPS: not disclosed \\
长飞光纤 601869 & 第三方一致预期页面 & 292.37 & -46.2\% & 15\% & B- & 收入/利润/EPS: not disclosed \\
烽火通信 600498 & 检索缺口 & n.d. & n.a. & 0\% & C & 收入/利润/EPS: not disclosed \\
中兴通讯 000063 & 第三方一致预期页面 & 43.69 & +22.5\% & 15\% & B & 收入/利润/EPS: not disclosed \\
仕佳光子 688313 & 第三方一致预期页面 & 157.53 & -12.3\% & 15\% & B- & 收入/利润/EPS: not disclosed \\
长光华芯 688048 & 第三方一致预期页面 & 86.84 & -78.1\% & 15\% & C+ & 收入/利润/EPS: not disclosed \\
光库科技 300620 & 第三方一致预期页面 & 177.00 & -51.2\% & 15\% & C+ & 收入/利润/EPS: not disclosed \\
通鼎互联 002491 & 第三方一致预期页面 & n.d. & n.a. & 0\% & C & 收入/利润/EPS: not disclosed \\
永鼎股份 600105 & 第三方一致预期页面 & n.d. & n.a. & 0\% & C & 收入/利润/EPS: not disclosed \\
长芯博创 300548 & 第三方一致预期页面 & n.d. & n.a. & 0\% & C+ & 收入/利润/EPS: not disclosed \\
德科立 688205 & 第三方一致预期页面 & n.d. & n.a. & 0\% & C+ & 收入/利润/EPS: not disclosed \\
腾景科技 688195 & 第三方一致预期页面 & 222.42 & +18.6\% & 15\% & B- & 收入/利润/EPS: not disclosed \\
联特科技 301205 & 媒体聚合摘要 & 271.98 & -13.9\% & 15\% & C+ & 收入/利润/EPS: not disclosed \\
兆龙互连 300913 & 第三方一致预期页面 & n.d. & n.a. & 0\% & C & 收入/利润/EPS: not disclosed \\
意华股份 002897 & 第三方一致预期页面 & n.d. & n.a. & 0\% & C & 收入/利润/EPS: not disclosed \\
神宇股份 300563 & 第三方一致预期页面 & n.d. & n.a. & 0\% & C & 收入/利润/EPS: not disclosed \\
杭电股份 603618 & 检索缺口 & n.d. & n.a. & 0\% & C & 收入/利润/EPS: not disclosed \\
\bottomrule
\end{tabularx}
\end{exhibitbox}
\sourcenote{券商/Street 一致预期对照包：data/broker\_street\_consensus\_20260626.md/json（含发布降级说明）。}

\section{情景解释}
中际旭创和新易盛的综合空间接近持平甚至略正，是因为 2026Q1 已体现强利润分母、TTM run-rate 扎实，且现价隐含增速（约 5\%--7\% 的 3 年归母 CAGR）其实低于 AStock 假设增速，说明并非无条件贵；但它们也不是无约束买入，核心风险是海外云厂商订单与 ASP 持续性。长飞光纤、亨通光电、中天科技、通鼎互联、永鼎股份、中兴通讯和烽火通信提供完整产业链的慢周期底盘，估值既要锚定现金流和运营商/网络项目节奏，也要承认成交额和市场共识给出的情绪溢价。源杰科技、长光华芯、仕佳光子、光库科技拥有更高战略稀缺性，但现价通常把长期技术期权前置，其中长光华芯为近零 EPS，只作观察池。光迅科技、华工科技、剑桥科技、太辰光的问题不是没有产业位置，而是利润分母、业务纯度或现金流还不足以支撑过高倍数。

\begin{exhibitbox}[表：估值假设桥 / Method and Assumption Bridge]
\centering\scriptsize
\begin{tabularx}{\exhibitboxwidth}{>{\bfseries\raggedright\arraybackslash}p{1.5cm} >{\centering\arraybackslash}p{1.0cm} >{\raggedright\arraybackslash\sloppy\hspace{0pt}}p{2.2cm} >{\raggedright\arraybackslash\sloppy\hspace{0pt}}X}
\toprule
\textbf{标的} & \textbf{方法} & \textbf{权重} & \textbf{二级校验与倍数理由} \\
\midrule
中际旭创 & PE/PEG & PE 100\% & 客户/订单持续性、毛利率保持、PEG 排名和 PS 交叉校验；800G/1.6T 规模和盈利兑现最强，但当前估值已经反映海外 AI 需求持续高景气。 \\
新易盛 & PE/PEG & PE 100\% & 客户/订单持续性、毛利率保持、PEG 排名和 PS 交叉校验；高增长光模块弹性强，2025--2026 年交付较强，但估值对单一客户和 ASP 假设敏感。 \\
天孚通信 & PE/PEG & PE 100\% & 客户/订单持续性、毛利率保持、PEG 排名和 PS 交叉校验；高毛利精密器件平台价值明确，但增速相对模块龙头放缓。 \\
光迅科技 & PE/PB/PS & PE 70\% / PB 10\% / PS 20\% & 数通收入占比、毛利率桥、收入规模和现金转化；产品线和央企平台优势存在，但盈利能力显著低于 AI 光模块龙头。 \\
华工科技 & PE/PB/PS & PE 70\% / PB 10\% / PS 20\% & 数通收入占比、毛利率桥、收入规模和现金转化；混合业务提供盈利稳定性，但削弱纯 AI 光模块弹性。 \\
剑桥科技 & PE/PB/PS & PE 45\% / PB 15\% / PS 40\% & 客户认证、收入规模、毛利率持续性和战略稀缺性；反转弹性存在，但 Q1 利润分母仍偏小，难以支撑过高目标价。 \\
太辰光 & PE/PB/PS & PE 45\% / PB 30\% / PS 25\% & 周期正常化 EPS、经营现金流、运营商/项目订单、光纤光缆价格；无源器件敞口真实，但 2026Q1 收入/利润下滑，估值需要交付修复。 \\
源杰科技 & PE/PB/PS & PE 45\% / PB 15\% / PS 40\% & 客户认证、收入规模、毛利率持续性和战略稀缺性；上游战略属性最强之一，但当前价格已经前置长期兑现路径。 \\
亨通光电 & PE/PB/PS & PE 45\% / PB 30\% / PS 25\% & 周期正常化 EPS、经营现金流、运营商/项目订单、光纤光缆价格；光纤光缆龙头兼具海缆和电网敞口，但 AI 数据中心纯度低于模块龙头。 \\
中天科技 & PE/PB/PS & PE 45\% / PB 30\% / PS 25\% & 周期正常化 EPS、经营现金流、运营商/项目订单、光纤光缆价格；电力海洋与通信线缆平台提供盈利韧性，但不能按纯 AI 光模块估值。 \\
长飞光纤 & PE/PB/PS & PE 45\% / PB 30\% / PS 25\% & 周期正常化 EPS、经营现金流、运营商/项目订单、光纤光缆价格；预制棒/光纤/光缆一体化稀缺性较强，但数据中心需求需体现在利用率上。 \\
烽火通信 & PE/PB/PS & PE 55\% / PB 20\% / PS 25\% & 在手订单、运营商资本开支、业务结构、现金流和非光通信业务稀释；运营商光网络平台具备战略意义，但 Q1 利润分母薄，周期慢于云侧 AI capex。 \\
中兴通讯 & PE/PB/PS & PE 55\% / PB 20\% / PS 25\% & 在手订单、运营商资本开支、业务结构、现金流和非光通信业务稀释；覆盖运营商、云网络和应用层，但光通信只是业务组合的一部分。 \\
仕佳光子 & PE/PB/PS & PE 45\% / PB 15\% / PS 40\% & 客户认证、收入规模、毛利率持续性和战略稀缺性；无源光芯片与 WDM/AWG/硅光生态相关性强，但估值需要高利用率支撑。 \\
长光华芯 & PE/PB/PS & PE 45\% / PB 15\% / PS 40\% & 客户认证、收入规模、毛利率持续性和战略稀缺性；激光芯片期权高，但 Q1 利润分母仍小，基准情景要求快速认证和放量。 \\
光库科技 & PE/PB/PS & PE 45\% / PB 30\% / PS 25\% & 周期正常化 EPS、经营现金流、运营商/项目订单、光纤光缆价格；光器件和薄膜铌酸锂期权提高战略相关性，但股价已经计入较长兑现路径。 \\
通鼎互联 & PE/PB/PS & PE 20\% / PB 30\% / PS 50\% & 现金转化、运营商/云侧线缆订单、线缆 ASP、AI/数据中心敞口纯度；属于完整产业链底座，但盈利质量和 AI 数据中心纯度低于模块龙头。 \\
永鼎股份 & PE/PB/PS & PE 20\% / PB 30\% / PS 50\% & 现金转化、运营商/云侧线缆订单、线缆 ASP、AI/数据中心敞口纯度；线缆和电信基础设施修复可见，但业务慢周期、纯度低于 AI 光模块。 \\
长芯博创 & PE/PB/PS & PE 70\% / PB 10\% / PS 20\% & 数通收入占比、毛利率桥、收入规模和现金转化；光模块/器件平台具备数通敞口，但当前估值要求高速产品持续交付。 \\
德科立 & PE/PB/PS & PE 45\% / PB 15\% / PS 40\% & 客户认证、收入规模、毛利率持续性和战略稀缺性；补足 DCI 和运营商相干传输链条，但 Q1 利润基数仍偏小。 \\
腾景科技 & PE/PB/PS & PE 45\% / PB 15\% / PS 40\% & 客户认证、收入规模、毛利率持续性和战略稀缺性；精密光学与高速光系统相关，但规模小，估值需要持续订单转化。 \\
联特科技 & PE/PB/PS & PE 45\% / PB 15\% / PS 40\% & 客户认证、收入规模、毛利率持续性和战略稀缺性；高速模块弹性大，但 Q1 利润分母很小，目标价高度依赖执行。 \\
兆龙互连 & PE/PB/PS & PE 40\% / PB 25\% / PS 35\% & AI 互连敞口纯度、产品结构、营运资本和毛利率稳定性；高速线缆/铜互连补足机柜内和企业网络层，但不是纯光模块替代品。 \\
意华股份 & PE/PB/PS & PE 40\% / PB 25\% / PS 35\% & AI 互连敞口纯度、产品结构、营运资本和毛利率稳定性；连接器敞口属于数据中心互连链，但混合业务需要纯度折价。 \\
神宇股份 & PE/PB/PS & PE 40\% / PB 25\% / PS 35\% & AI 互连敞口纯度、产品结构、营运资本和毛利率稳定性；线缆/互连拓宽产业链覆盖，但光通信纯度和 Q1 利润规模有限。 \\
杭电股份 & PE/PB/PS & PE 45\% / PB 30\% / PS 25\% & 周期正常化 EPS、经营现金流、运营商/项目订单、光纤光缆价格；AI 驱动光纤需求带动 H1 2026 从 2025 年亏损大幅扭亏，但基础业务为低毛利线缆，复苏需在一次预告之外持续确认。 \\
\bottomrule
\end{tabularx}
\end{exhibitbox}

\section{单票估值卡与压力测试}
本节把最终估值表拆成单票估值卡。每张卡都明确估值方法、EPS/收入/BVPS 分母、PE/PB/PS 组件、三档价值和动作，方便后续在行情变化后直接替换价格、财务分母或估值权重重算。

\subsection{中际旭创（300308）估值卡}
\begin{exhibitbox}[估值卡：中际旭创]
\centering
\scriptsize
\begin{tabularx}{\exhibitboxwidth}{>{\bfseries\raggedright\arraybackslash}p{2.5cm} >{\raggedleft\arraybackslash}p{1.6cm} >{\raggedleft\arraybackslash}p{1.6cm} >{\raggedleft\arraybackslash}p{1.6cm} >{\raggedright\arraybackslash\sloppy\hspace{0pt}}X}
\toprule
\textbf{项目} & \textbf{熊市} & \textbf{基准} & \textbf{牛市} & \textbf{说明} \\
\midrule
估值方法 & \multicolumn{3}{r}{PE/PEG} & 基于正常化 EPS 的 forward PE/PEG，用于盈利兑现强且订单持续的 AI 光模块龙头；权重 PE 100\% \\
EPS 分母 & 22.52 & 29.73 & 36.27 & 2026E 由 Q1 季节性折算，2027/2028E 由公司层面成长假设外推 \\
收入/股 & 76.57 & 101.07 & 123.31 & PS 组件用于校验收入规模、业务纯度和主题期权，不替代交付验证 \\
BVPS & 32.03 & 32.03 & 32.03 & PB 组件用于资产、项目和混合业务的底部校验 \\
PE 组件 & 788.26 & 1337.79 & 1781.07 & PE 档位 35x/45x/55x，只按权重进入综合价值 \\
PB 组件 & n.a. & n.a. & n.a. & PB 档位 0.0x/0.0x/0.0x，非资产型公司为辅助或 n.a. \\
PS 组件 & n.a. & n.a. & n.a. & PS 档位 0.0x/0.0x/0.0x，牛市价值按 12\% 折现 \\
内在价值锚 & 788.26 & 1337.79 & 1781.07 & 来自业务模型组件加权，不是所有标的统一套 PE \\
市场情绪锚 & \multicolumn{3}{r}{1003.11} & 成交额分位 +100.0\%，情绪状态 强共识/高拥挤，当前价较内在锚溢价 -6.3\% \\
综合目标价 & \multicolumn{3}{r}{1229.56} & 内在/市场/券商权重：内在65\% / 市场20\% / 券商15\% \\
当前价格 & \multicolumn{3}{r}{1253.89} & 2026-06-26 收盘后行情包 \\
空间与动作 & \multicolumn{3}{r}{-1.9\%；\stance{riskamber}{中性}} & 证据质量 A-；动作由综合目标价相对当前价决定。 \\

\bottomrule
\end{tabularx}
\end{exhibitbox}

压力测试上，如果 2027E EPS 下修 15\%，仅 PE 组件受影响，中际旭创 的内在价值锚降至约 CNY1137.12；若市场风险偏好下行导致所有估值组件倍数降 20\%，内在锚降至约 CNY1070.23。两者同时发生时，内在锚约为 CNY909.70。若成交额分位回落、题材拥挤退潮或券商/市场共识下修，市场情绪锚也会同步下移，综合目标价会向内在价值锚收敛。这就是本报告要求用 Q2/Q3 交付、现金流和订单质量确认情绪溢价的原因。


\subsection{新易盛（300502）估值卡}
\begin{exhibitbox}[估值卡：新易盛]
\centering
\scriptsize
\begin{tabularx}{\exhibitboxwidth}{>{\bfseries\raggedright\arraybackslash}p{2.5cm} >{\raggedleft\arraybackslash}p{1.6cm} >{\raggedleft\arraybackslash}p{1.6cm} >{\raggedleft\arraybackslash}p{1.6cm} >{\raggedright\arraybackslash\sloppy\hspace{0pt}}X}
\toprule
\textbf{项目} & \textbf{熊市} & \textbf{基准} & \textbf{牛市} & \textbf{说明} \\
\midrule
估值方法 & \multicolumn{3}{r}{PE/PEG} & 基于正常化 EPS 的 forward PE/PEG，用于盈利兑现强且订单持续的 AI 光模块龙头；权重 PE 100\% \\
EPS 分母 & 12.17 & 15.83 & 19.31 & 2026E 由 Q1 季节性折算，2027/2028E 由公司层面成长假设外推 \\
收入/股 & 36.51 & 47.46 & 57.90 & PS 组件用于校验收入规模、业务纯度和主题期权，不替代交付验证 \\
BVPS & 20.50 & 20.50 & 20.50 & PB 组件用于资产、项目和混合业务的底部校验 \\
PE 组件 & 389.57 & 633.04 & 861.96 & PE 档位 32x/40x/50x，只按权重进入综合价值 \\
PB 组件 & n.a. & n.a. & n.a. & PB 档位 0.0x/0.0x/0.0x，非资产型公司为辅助或 n.a. \\
PS 组件 & n.a. & n.a. & n.a. & PS 档位 0.0x/0.0x/0.0x，牛市价值按 12\% 折现 \\
内在价值锚 & 389.57 & 633.04 & 861.96 & 来自业务模型组件加权，不是所有标的统一套 PE \\
市场情绪锚 & \multicolumn{3}{r}{452.80} & 成交额分位 +96.2\%，情绪状态 强共识/高拥挤，当前价较内在锚溢价 -10.6\% \\
综合目标价 & \multicolumn{3}{r}{576.55} & 内在/市场/券商权重：内在65\% / 市场20\% / 券商15\% \\
当前价格 & \multicolumn{3}{r}{566.00} & 2026-06-26 收盘后行情包 \\
空间与动作 & \multicolumn{3}{r}{+1.9\%；\stance{riskamber}{中性}} & 证据质量 A-；动作由综合目标价相对当前价决定。 \\

\bottomrule
\end{tabularx}
\end{exhibitbox}

压力测试上，如果 2027E EPS 下修 15\%，仅 PE 组件受影响，新易盛 的内在价值锚降至约 CNY538.09；若市场风险偏好下行导致所有估值组件倍数降 20\%，内在锚降至约 CNY506.43。两者同时发生时，内在锚约为 CNY430.47。若成交额分位回落、题材拥挤退潮或券商/市场共识下修，市场情绪锚也会同步下移，综合目标价会向内在价值锚收敛。这就是本报告要求用 Q2/Q3 交付、现金流和订单质量确认情绪溢价的原因。


\subsection{天孚通信（300394）估值卡}
\begin{exhibitbox}[估值卡：天孚通信]
\centering
\scriptsize
\begin{tabularx}{\exhibitboxwidth}{>{\bfseries\raggedright\arraybackslash}p{2.5cm} >{\raggedleft\arraybackslash}p{1.6cm} >{\raggedleft\arraybackslash}p{1.6cm} >{\raggedleft\arraybackslash}p{1.6cm} >{\raggedright\arraybackslash\sloppy\hspace{0pt}}X}
\toprule
\textbf{项目} & \textbf{熊市} & \textbf{基准} & \textbf{牛市} & \textbf{说明} \\
\midrule
估值方法 & \multicolumn{3}{r}{PE/PEG} & 基于正常化 EPS 的 forward PE/PEG，用于盈利兑现强且订单持续的 AI 光模块龙头；权重 PE 100\% \\
EPS 分母 & 2.64 & 3.27 & 3.86 & 2026E 由 Q1 季节性折算，2027/2028E 由公司层面成长假设外推 \\
收入/股 & 7.13 & 8.84 & 10.43 & PS 组件用于校验收入规模、业务纯度和主题期权，不替代交付验证 \\
BVPS & 7.69 & 7.69 & 7.69 & PB 组件用于资产、项目和混合业务的底部校验 \\
PE 组件 & 92.33 & 147.20 & 189.54 & PE 档位 35x/45x/55x，只按权重进入综合价值 \\
PB 组件 & n.a. & n.a. & n.a. & PB 档位 0.0x/0.0x/0.0x，非资产型公司为辅助或 n.a. \\
PS 组件 & n.a. & n.a. & n.a. & PS 档位 0.0x/0.0x/0.0x，牛市价值按 12\% 折现 \\
内在价值锚 & 92.33 & 147.20 & 189.54 & 来自业务模型组件加权，不是所有标的统一套 PE \\
市场情绪锚 & \multicolumn{3}{r}{254.40} & 成交额分位 +84.6\%，情绪状态 强共识/高拥挤，当前价较内在锚溢价 +116.0\% \\
综合目标价 & \multicolumn{3}{r}{254.40} & 内在/市场/券商权重：内在57\% / 市场28\% / 券商15\% \\
当前价格 & \multicolumn{3}{r}{318.00} & 2026-06-26 收盘后行情包 \\
空间与动作 & \multicolumn{3}{r}{-20.0\%；\stance{riskamber}{中性观察}} & 证据质量 B+；市场高成交已给出情绪溢价，维持中性观察并等待利润/现金流追认。 \\

\bottomrule
\end{tabularx}
\end{exhibitbox}

压力测试上，如果 2027E EPS 下修 15\%，仅 PE 组件受影响，天孚通信 的内在价值锚降至约 CNY125.12；若市场风险偏好下行导致所有估值组件倍数降 20\%，内在锚降至约 CNY117.76。两者同时发生时，内在锚约为 CNY100.09。若成交额分位回落、题材拥挤退潮或券商/市场共识下修，市场情绪锚也会同步下移，综合目标价会向内在价值锚收敛。这就是本报告要求用 Q2/Q3 交付、现金流和订单质量确认情绪溢价的原因。


\subsection{光迅科技（002281）估值卡}
\begin{exhibitbox}[估值卡：光迅科技]
\centering
\scriptsize
\begin{tabularx}{\exhibitboxwidth}{>{\bfseries\raggedright\arraybackslash}p{2.5cm} >{\raggedleft\arraybackslash}p{1.6cm} >{\raggedleft\arraybackslash}p{1.6cm} >{\raggedleft\arraybackslash}p{1.6cm} >{\raggedright\arraybackslash\sloppy\hspace{0pt}}X}
\toprule
\textbf{项目} & \textbf{熊市} & \textbf{基准} & \textbf{牛市} & \textbf{说明} \\
\midrule
估值方法 & \multicolumn{3}{r}{PE/PB/PS} & 70\% PE + 10\% PB + 20\% PS，用于有正 EPS 但 AI 纯度不完美的模块/器件平台；权重 PE 70\% / PB 10\% / PS 20\% \\
EPS 分母 & 1.36 & 1.64 & 1.88 & 2026E 由 Q1 季节性折算，2027/2028E 由公司层面成长假设外推 \\
收入/股 & 15.76 & 18.91 & 21.75 & PS 组件用于校验收入规模、业务纯度和主题期权，不替代交付验证 \\
BVPS & 12.82 & 12.82 & 12.82 & PB 组件用于资产、项目和混合业务的底部校验 \\
PE 组件 & 38.18 & 57.27 & 75.61 & PE 档位 28x/35x/45x，只按权重进入综合价值 \\
PB 组件 & 32.05 & 51.27 & 57.23 & PB 档位 2.5x/4.0x/5.0x，非资产型公司为辅助或 n.a. \\
PS 组件 & 31.52 & 66.20 & 97.11 & PS 档位 2.0x/3.5x/5.0x，牛市价值按 12\% 折现 \\
内在价值锚 & 36.24 & 58.46 & 78.07 & 来自业务模型组件加权，不是所有标的统一套 PE \\
市场情绪锚 & \multicolumn{3}{r}{187.47} & 成交额分位 +80.8\%，情绪状态 强共识/高拥挤，当前价较内在锚溢价 +311.1\% \\
综合目标价 & \multicolumn{3}{r}{187.47} & 内在/市场/券商权重：内在47\% / 市场38\% / 券商15\% \\
当前价格 & \multicolumn{3}{r}{240.34} & 2026-06-26 收盘后行情包 \\
空间与动作 & \multicolumn{3}{r}{-22.0\%；\stance{riskamber}{中性观察}} & 证据质量 B；市场高成交已给出情绪溢价，维持中性观察并等待利润/现金流追认。 \\

\bottomrule
\end{tabularx}
\end{exhibitbox}

压力测试上，如果 2027E EPS 下修 15\%，仅 PE 组件受影响，光迅科技 的内在价值锚降至约 CNY52.44；若市场风险偏好下行导致所有估值组件倍数降 20\%，内在锚降至约 CNY46.77。两者同时发生时，内在锚约为 CNY41.96。若成交额分位回落、题材拥挤退潮或券商/市场共识下修，市场情绪锚也会同步下移，综合目标价会向内在价值锚收敛。这就是本报告要求用 Q2/Q3 交付、现金流和订单质量确认情绪溢价的原因。


\subsection{华工科技（000988）估值卡}
\begin{exhibitbox}[估值卡：华工科技]
\centering
\scriptsize
\begin{tabularx}{\exhibitboxwidth}{>{\bfseries\raggedright\arraybackslash}p{2.5cm} >{\raggedleft\arraybackslash}p{1.6cm} >{\raggedleft\arraybackslash}p{1.6cm} >{\raggedleft\arraybackslash}p{1.6cm} >{\raggedright\arraybackslash\sloppy\hspace{0pt}}X}
\toprule
\textbf{项目} & \textbf{熊市} & \textbf{基准} & \textbf{牛市} & \textbf{说明} \\
\midrule
估值方法 & \multicolumn{3}{r}{PE/PB/PS} & 70\% PE + 10\% PB + 20\% PS，用于有正 EPS 但 AI 纯度不完美的模块/器件平台；权重 PE 70\% / PB 10\% / PS 20\% \\
EPS 分母 & 2.56 & 3.02 & 3.44 & 2026E 由 Q1 季节性折算，2027/2028E 由公司层面成长假设外推 \\
收入/股 & 17.10 & 20.18 & 23.01 & PS 组件用于校验收入规模、业务纯度和主题期权，不替代交付验证 \\
BVPS & 11.62 & 11.62 & 11.62 & PB 组件用于资产、项目和混合业务的底部校验 \\
PE 组件 & 64.00 & 105.73 & 138.36 & PE 档位 25x/35x/45x，只按权重进入综合价值 \\
PB 组件 & 29.06 & 46.49 & 51.89 & PB 档位 2.5x/4.0x/5.0x，非资产型公司为辅助或 n.a. \\
PS 组件 & 34.21 & 70.64 & 102.72 & PS 档位 2.0x/3.5x/5.0x，牛市价值按 12\% 折现 \\
内在价值锚 & 54.55 & 92.79 & 122.59 & 来自业务模型组件加权，不是所有标的统一套 PE \\
市场情绪锚 & \multicolumn{3}{r}{128.75} & 成交额分位 +76.9\%，情绪状态 强共识/高拥挤，当前价较内在锚溢价 +73.5\% \\
综合目标价 & \multicolumn{3}{r}{128.75} & 内在/市场/券商权重：内在47\% / 市场38\% / 券商15\% \\
当前价格 & \multicolumn{3}{r}{160.94} & 2026-06-26 收盘后行情包 \\
空间与动作 & \multicolumn{3}{r}{-20.0\%；\stance{riskamber}{中性观察}} & 证据质量 B；市场高成交已给出情绪溢价，维持中性观察并等待利润/现金流追认。 \\

\bottomrule
\end{tabularx}
\end{exhibitbox}

压力测试上，如果 2027E EPS 下修 15\%，仅 PE 组件受影响，华工科技 的内在价值锚降至约 CNY81.69；若市场风险偏好下行导致所有估值组件倍数降 20\%，内在锚降至约 CNY74.23。两者同时发生时，内在锚约为 CNY65.35。若成交额分位回落、题材拥挤退潮或券商/市场共识下修，市场情绪锚也会同步下移，综合目标价会向内在价值锚收敛。这就是本报告要求用 Q2/Q3 交付、现金流和订单质量确认情绪溢价的原因。


\subsection{剑桥科技（603083）估值卡}
\begin{exhibitbox}[估值卡：剑桥科技]
\centering
\scriptsize
\begin{tabularx}{\exhibitboxwidth}{>{\bfseries\raggedright\arraybackslash}p{2.5cm} >{\raggedleft\arraybackslash}p{1.6cm} >{\raggedleft\arraybackslash}p{1.6cm} >{\raggedleft\arraybackslash}p{1.6cm} >{\raggedright\arraybackslash\sloppy\hspace{0pt}}X}
\toprule
\textbf{项目} & \textbf{熊市} & \textbf{基准} & \textbf{牛市} & \textbf{说明} \\
\midrule
估值方法 & \multicolumn{3}{r}{PE/PB/PS} & 45\% PE + 15\% PB + 40\% PS，用于小规模但具战略稀缺性的光芯片、相干模块和精密光学；权重 PE 45\% / PB 15\% / PS 40\% \\
EPS 分母 & 1.70 & 2.21 & 2.76 & 2026E 由 Q1 季节性折算，2027/2028E 由公司层面成长假设外推 \\
收入/股 & 18.48 & 24.03 & 30.04 & PS 组件用于校验收入规模、业务纯度和主题期权，不替代交付验证 \\
BVPS & 21.45 & 21.45 & 21.45 & PB 组件用于资产、项目和混合业务的底部校验 \\
PE 组件 & 51.00 & 99.45 & 135.66 & PE 档位 30x/45x/55x，只按权重进入综合价值 \\
PB 组件 & 53.61 & 85.78 & 105.31 & PB 档位 2.5x/4.0x/5.5x，非资产型公司为辅助或 n.a. \\
PS 组件 & 55.45 & 120.14 & 187.72 & PS 档位 3.0x/5.0x/7.0x，牛市价值按 12\% 折现 \\
内在价值锚 & 53.17 & 105.68 & 151.93 & 来自业务模型组件加权，不是所有标的统一套 PE \\
市场情绪锚 & \multicolumn{3}{r}{177.06} & 成交额分位 +42.3\%，情绪状态 活跃共识，当前价较内在锚溢价 +132.7\% \\
综合目标价 & \multicolumn{3}{r}{177.06} & 内在/市场/券商权重：内在47\% / 市场38\% / 券商15\% \\
当前价格 & \multicolumn{3}{r}{245.92} & 2026-06-26 收盘后行情包 \\
空间与动作 & \multicolumn{3}{r}{-28.0\%；\stance{riskamber}{中性观察}} & 证据质量 C+；市场高成交已给出情绪溢价，维持中性观察并等待利润/现金流追认。 \\

\bottomrule
\end{tabularx}
\end{exhibitbox}

压力测试上，如果 2027E EPS 下修 15\%，仅 PE 组件受影响，剑桥科技 的内在价值锚降至约 CNY98.96；若市场风险偏好下行导致所有估值组件倍数降 20\%，内在锚降至约 CNY84.54。两者同时发生时，内在锚约为 CNY79.17。若成交额分位回落、题材拥挤退潮或券商/市场共识下修，市场情绪锚也会同步下移，综合目标价会向内在价值锚收敛。这就是本报告要求用 Q2/Q3 交付、现金流和订单质量确认情绪溢价的原因。


\subsection{太辰光（300570）估值卡}
\begin{exhibitbox}[估值卡：太辰光]
\centering
\scriptsize
\begin{tabularx}{\exhibitboxwidth}{>{\bfseries\raggedright\arraybackslash}p{2.5cm} >{\raggedleft\arraybackslash}p{1.6cm} >{\raggedleft\arraybackslash}p{1.6cm} >{\raggedleft\arraybackslash}p{1.6cm} >{\raggedright\arraybackslash\sloppy\hspace{0pt}}X}
\toprule
\textbf{项目} & \textbf{熊市} & \textbf{基准} & \textbf{牛市} & \textbf{说明} \\
\midrule
估值方法 & \multicolumn{3}{r}{PE/PB/PS} & 45\% 周期正常化 PE + 30\% PB/ROE + 25\% PS，用于资产较重的光纤光缆和海缆平台；权重 PE 45\% / PB 30\% / PS 25\% \\
EPS 分母 & 1.16 & 1.33 & 1.49 & 2026E 由 Q1 季节性折算，2027/2028E 由公司层面成长假设外推 \\
收入/股 & 6.01 & 6.91 & 7.74 & PS 组件用于校验收入规模、业务纯度和主题期权，不替代交付验证 \\
BVPS & 7.79 & 7.79 & 7.79 & PB 组件用于资产、项目和混合业务的底部校验 \\
PE 组件 & 29.00 & 46.69 & 60.03 & PE 档位 25x/35x/45x，只按权重进入综合价值 \\
PB 组件 & 10.13 & 14.02 & 16.69 & PB 档位 1.3x/1.8x/2.4x，非资产型公司为辅助或 n.a. \\
PS 组件 & 3.31 & 5.53 & 7.26 & PS 档位 0.55x/0.8x/1.05x，牛市价值按 12\% 折现 \\
内在价值锚 & 16.91 & 26.60 & 33.83 & 来自业务模型组件加权，不是所有标的统一套 PE \\
市场情绪锚 & \multicolumn{3}{r}{163.87} & 成交额分位 +53.8\%，情绪状态 活跃共识，当前价较内在锚溢价 +780.1\% \\
综合目标价 & \multicolumn{3}{r}{163.87} & 内在/市场/券商权重：内在38\% / 市场47\% / 券商15\% \\
当前价格 & \multicolumn{3}{r}{234.10} & 2026-06-26 收盘后行情包 \\
空间与动作 & \multicolumn{3}{r}{-30.0\%；\stance{riskamber}{中性观察}} & 证据质量 C；市场高成交已给出情绪溢价，维持中性观察并等待利润/现金流追认。 \\

\bottomrule
\end{tabularx}
\end{exhibitbox}

压力测试上，如果 2027E EPS 下修 15\%，仅 PE 组件受影响，太辰光 的内在价值锚降至约 CNY23.45；若市场风险偏好下行导致所有估值组件倍数降 20\%，内在锚降至约 CNY21.28。两者同时发生时，内在锚约为 CNY18.76。若成交额分位回落、题材拥挤退潮或券商/市场共识下修，市场情绪锚也会同步下移，综合目标价会向内在价值锚收敛。这就是本报告要求用 Q2/Q3 交付、现金流和订单质量确认情绪溢价的原因。


\subsection{源杰科技（688498）估值卡}
\begin{exhibitbox}[估值卡：源杰科技]
\centering
\scriptsize
\begin{tabularx}{\exhibitboxwidth}{>{\bfseries\raggedright\arraybackslash}p{2.5cm} >{\raggedleft\arraybackslash}p{1.6cm} >{\raggedleft\arraybackslash}p{1.6cm} >{\raggedleft\arraybackslash}p{1.6cm} >{\raggedright\arraybackslash\sloppy\hspace{0pt}}X}
\toprule
\textbf{项目} & \textbf{熊市} & \textbf{基准} & \textbf{牛市} & \textbf{说明} \\
\midrule
估值方法 & \multicolumn{3}{r}{PE/PB/PS} & 45\% PE + 15\% PB + 40\% PS，用于小规模但具战略稀缺性的光芯片、相干模块和精密光学；权重 PE 45\% / PB 15\% / PS 40\% \\
EPS 分母 & 9.09 & 12.54 & 16.30 & 2026E 由 Q1 季节性折算，2027/2028E 由公司层面成长假设外推 \\
收入/股 & 17.99 & 24.83 & 32.28 & PS 组件用于校验收入规模、业务纯度和主题期权，不替代交付验证 \\
BVPS & 29.38 & 29.38 & 29.38 & PB 组件用于资产、项目和混合业务的底部校验 \\
PE 组件 & 363.48 & 815.10 & 1237.21 & PE 档位 40x/65x/85x，只按权重进入综合价值 \\
PB 组件 & 73.44 & 117.50 & 144.25 & PB 档位 2.5x/4.0x/5.5x，非资产型公司为辅助或 n.a. \\
PS 组件 & 53.97 & 124.14 & 201.73 & PS 档位 3.0x/5.0x/7.0x，牛市价值按 12\% 折现 \\
内在价值锚 & 196.17 & 434.08 & 659.07 & 来自业务模型组件加权，不是所有标的统一套 PE \\
市场情绪锚 & \multicolumn{3}{r}{1421.53} & 成交额分位 +65.4\%，情绪状态 强共识/高拥挤，当前价较内在锚溢价 +309.4\% \\
综合目标价 & \multicolumn{3}{r}{1421.53} & 内在/市场/券商权重：内在47\% / 市场38\% / 券商15\% \\
当前价格 & \multicolumn{3}{r}{1776.91} & 2026-06-26 收盘后行情包 \\
空间与动作 & \multicolumn{3}{r}{-20.0\%；\stance{riskamber}{中性观察}} & 证据质量 B；市场高成交已给出情绪溢价，维持中性观察并等待利润/现金流追认。 \\

\bottomrule
\end{tabularx}
\end{exhibitbox}

压力测试上，如果 2027E EPS 下修 15\%，仅 PE 组件受影响，源杰科技 的内在价值锚降至约 CNY379.06；若市场风险偏好下行导致所有估值组件倍数降 20\%，内在锚降至约 CNY347.26。两者同时发生时，内在锚约为 CNY303.25。若成交额分位回落、题材拥挤退潮或券商/市场共识下修，市场情绪锚也会同步下移，综合目标价会向内在价值锚收敛。这就是本报告要求用 Q2/Q3 交付、现金流和订单质量确认情绪溢价的原因。


\subsection{亨通光电（600487）估值卡}
\begin{exhibitbox}[估值卡：亨通光电]
\centering
\scriptsize
\begin{tabularx}{\exhibitboxwidth}{>{\bfseries\raggedright\arraybackslash}p{2.5cm} >{\raggedleft\arraybackslash}p{1.6cm} >{\raggedleft\arraybackslash}p{1.6cm} >{\raggedleft\arraybackslash}p{1.6cm} >{\raggedright\arraybackslash\sloppy\hspace{0pt}}X}
\toprule
\textbf{项目} & \textbf{熊市} & \textbf{基准} & \textbf{牛市} & \textbf{说明} \\
\midrule
估值方法 & \multicolumn{3}{r}{PE/PB/PS} & 45\% 周期正常化 PE + 30\% PB/ROE + 25\% PS，用于资产较重的光纤光缆和海缆平台；权重 PE 45\% / PB 30\% / PS 25\% \\
EPS 分母 & 1.90 & 2.13 & 2.34 & 2026E 由 Q1 季节性折算，2027/2028E 由公司层面成长假设外推 \\
收入/股 & 30.60 & 34.27 & 37.70 & PS 组件用于校验收入规模、业务纯度和主题期权，不替代交付验证 \\
BVPS & 13.14 & 13.14 & 13.14 & PB 组件用于资产、项目和混合业务的底部校验 \\
PE 组件 & 34.22 & 51.11 & 62.74 & PE 档位 18x/24x/30x，只按权重进入综合价值 \\
PB 组件 & 17.08 & 23.65 & 28.16 & PB 档位 1.3x/1.8x/2.4x，非资产型公司为辅助或 n.a. \\
PS 组件 & 16.83 & 27.42 & 35.34 & PS 档位 0.55x/0.8x/1.05x，牛市价值按 12\% 折现 \\
内在价值锚 & 24.73 & 36.95 & 45.52 & 来自业务模型组件加权，不是所有标的统一套 PE \\
市场情绪锚 & \multicolumn{3}{r}{86.43} & 成交额分位 +88.5\%，情绪状态 强共识/高拥挤，当前价较内在锚溢价 +199.9\% \\
综合目标价 & \multicolumn{3}{r}{86.43} & 内在/市场/券商权重：内在38\% / 市场47\% / 券商15\% \\
当前价格 & \multicolumn{3}{r}{110.81} & 2026-06-26 收盘后行情包 \\
空间与动作 & \multicolumn{3}{r}{-22.0\%；\stance{riskamber}{中性观察}} & 证据质量 B；市场高成交已给出情绪溢价，维持中性观察并等待利润/现金流追认。 \\

\bottomrule
\end{tabularx}
\end{exhibitbox}

压力测试上，如果 2027E EPS 下修 15\%，仅 PE 组件受影响，亨通光电 的内在价值锚降至约 CNY33.50；若市场风险偏好下行导致所有估值组件倍数降 20\%，内在锚降至约 CNY29.56。两者同时发生时，内在锚约为 CNY26.80。若成交额分位回落、题材拥挤退潮或券商/市场共识下修，市场情绪锚也会同步下移，综合目标价会向内在价值锚收敛。这就是本报告要求用 Q2/Q3 交付、现金流和订单质量确认情绪溢价的原因。


\subsection{中天科技（600522）估值卡}
\begin{exhibitbox}[估值卡：中天科技]
\centering
\scriptsize
\begin{tabularx}{\exhibitboxwidth}{>{\bfseries\raggedright\arraybackslash}p{2.5cm} >{\raggedleft\arraybackslash}p{1.6cm} >{\raggedleft\arraybackslash}p{1.6cm} >{\raggedleft\arraybackslash}p{1.6cm} >{\raggedright\arraybackslash\sloppy\hspace{0pt}}X}
\toprule
\textbf{项目} & \textbf{熊市} & \textbf{基准} & \textbf{牛市} & \textbf{说明} \\
\midrule
估值方法 & \multicolumn{3}{r}{PE/PB/PS} & 45\% 周期正常化 PE + 30\% PB/ROE + 25\% PS，用于资产较重的光纤光缆和海缆平台；权重 PE 45\% / PB 30\% / PS 25\% \\
EPS 分母 & 1.13 & 1.27 & 1.40 & 2026E 由 Q1 季节性折算，2027/2028E 由公司层面成长假设外推 \\
收入/股 & 16.21 & 18.15 & 19.97 & PS 组件用于校验收入规模、业务纯度和主题期权，不替代交付验证 \\
BVPS & 11.14 & 11.14 & 11.14 & PB 组件用于资产、项目和混合业务的底部校验 \\
PE 组件 & 20.40 & 30.46 & 37.40 & PE 档位 18x/24x/30x，只按权重进入综合价值 \\
PB 组件 & 14.48 & 20.05 & 23.87 & PB 档位 1.3x/1.8x/2.4x，非资产型公司为辅助或 n.a. \\
PS 组件 & 8.91 & 14.52 & 18.72 & PS 档位 0.55x/0.8x/1.05x，牛市价值按 12\% 折现 \\
内在价值锚 & 15.75 & 23.35 & 28.67 & 来自业务模型组件加权，不是所有标的统一套 PE \\
市场情绪锚 & \multicolumn{3}{r}{47.81} & 成交额分位 +92.3\%，情绪状态 强共识/高拥挤，当前价较内在锚溢价 +162.4\% \\
综合目标价 & \multicolumn{3}{r}{47.81} & 内在/市场/券商权重：内在38\% / 市场47\% / 券商15\% \\
当前价格 & \multicolumn{3}{r}{61.29} & 2026-06-26 收盘后行情包 \\
空间与动作 & \multicolumn{3}{r}{-22.0\%；\stance{riskamber}{中性观察}} & 证据质量 B；市场高成交已给出情绪溢价，维持中性观察并等待利润/现金流追认。 \\

\bottomrule
\end{tabularx}
\end{exhibitbox}

压力测试上，如果 2027E EPS 下修 15\%，仅 PE 组件受影响，中天科技 的内在价值锚降至约 CNY21.30；若市场风险偏好下行导致所有估值组件倍数降 20\%，内在锚降至约 CNY18.68。两者同时发生时，内在锚约为 CNY17.04。若成交额分位回落、题材拥挤退潮或券商/市场共识下修，市场情绪锚也会同步下移，综合目标价会向内在价值锚收敛。这就是本报告要求用 Q2/Q3 交付、现金流和订单质量确认情绪溢价的原因。


\subsection{长飞光纤（601869）估值卡}
\begin{exhibitbox}[估值卡：长飞光纤]
\centering
\scriptsize
\begin{tabularx}{\exhibitboxwidth}{>{\bfseries\raggedright\arraybackslash}p{2.5cm} >{\raggedleft\arraybackslash}p{1.6cm} >{\raggedleft\arraybackslash}p{1.6cm} >{\raggedleft\arraybackslash}p{1.6cm} >{\raggedright\arraybackslash\sloppy\hspace{0pt}}X}
\toprule
\textbf{项目} & \textbf{熊市} & \textbf{基准} & \textbf{牛市} & \textbf{说明} \\
\midrule
估值方法 & \multicolumn{3}{r}{PE/PB/PS} & 45\% 周期正常化 PE + 30\% PB/ROE + 25\% PS，用于资产较重的光纤光缆和海缆平台；权重 PE 45\% / PB 30\% / PS 25\% \\
EPS 分母 & 2.61 & 3.03 & 3.39 & 2026E 由 Q1 季节性折算，2027/2028E 由公司层面成长假设外推 \\
收入/股 & 19.47 & 22.58 & 25.30 & PS 组件用于校验收入规模、业务纯度和主题期权，不替代交付验证 \\
BVPS & 17.19 & 17.19 & 17.19 & PB 组件用于资产、项目和混合业务的底部校验 \\
PE 组件 & 52.17 & 84.73 & 105.91 & PE 档位 20x/28x/35x，只按权重进入综合价值 \\
PB 组件 & 22.35 & 30.94 & 36.84 & PB 档位 1.3x/1.8x/2.4x，非资产型公司为辅助或 n.a. \\
PS 组件 & 10.71 & 18.07 & 23.71 & PS 档位 0.55x/0.8x/1.05x，牛市价值按 12\% 折现 \\
内在价值锚 & 32.86 & 51.93 & 64.64 & 来自业务模型组件加权，不是所有标的统一套 PE \\
市场情绪锚 & \multicolumn{3}{r}{434.40} & 成交额分位 +61.5\%，情绪状态 强共识/高拥挤，当前价较内在锚溢价 +945.7\% \\
综合目标价 & \multicolumn{3}{r}{434.40} & 内在/市场/券商权重：内在38\% / 市场47\% / 券商15\% \\
当前价格 & \multicolumn{3}{r}{543.00} & 2026-06-26 收盘后行情包 \\
空间与动作 & \multicolumn{3}{r}{-20.0\%；\stance{riskamber}{中性观察}} & 证据质量 B+；市场高成交已给出情绪溢价，维持中性观察并等待利润/现金流追认。 \\

\bottomrule
\end{tabularx}
\end{exhibitbox}

压力测试上，如果 2027E EPS 下修 15\%，仅 PE 组件受影响，长飞光纤 的内在价值锚降至约 CNY46.21；若市场风险偏好下行导致所有估值组件倍数降 20\%，内在锚降至约 CNY41.54。两者同时发生时，内在锚约为 CNY36.97。若成交额分位回落、题材拥挤退潮或券商/市场共识下修，市场情绪锚也会同步下移，综合目标价会向内在价值锚收敛。这就是本报告要求用 Q2/Q3 交付、现金流和订单质量确认情绪溢价的原因。


\subsection{烽火通信（600498）估值卡}
\begin{exhibitbox}[估值卡：烽火通信]
\centering
\scriptsize
\begin{tabularx}{\exhibitboxwidth}{>{\bfseries\raggedright\arraybackslash}p{2.5cm} >{\raggedleft\arraybackslash}p{1.6cm} >{\raggedleft\arraybackslash}p{1.6cm} >{\raggedleft\arraybackslash}p{1.6cm} >{\raggedright\arraybackslash\sloppy\hspace{0pt}}X}
\toprule
\textbf{项目} & \textbf{熊市} & \textbf{基准} & \textbf{牛市} & \textbf{说明} \\
\midrule
估值方法 & \multicolumn{3}{r}{PE/PB/PS} & 55\% PE + 20\% PB + 25\% PS，用于运营商/云网络设备业务；权重 PE 55\% / PB 20\% / PS 25\% \\
EPS 分母 & 0.14 & 0.16 & 0.18 & 2026E 由 Q1 季节性折算，2027/2028E 由公司层面成长假设外推 \\
收入/股 & 16.02 & 18.42 & 20.63 & PS 组件用于校验收入规模、业务纯度和主题期权，不替代交付验证 \\
BVPS & 12.94 & 12.94 & 12.94 & PB 组件用于资产、项目和混合业务的底部校验 \\
PE 组件 & 2.73 & 4.39 & 5.49 & PE 档位 20x/28x/35x，只按权重进入综合价值 \\
PB 组件 & 18.11 & 27.17 & 32.35 & PB 档位 1.4x/2.1x/2.8x，非资产型公司为辅助或 n.a. \\
PS 组件 & 10.41 & 18.42 & 24.87 & PS 档位 0.65x/1.0x/1.35x，牛市价值按 12\% 折现 \\
内在价值锚 & 7.73 & 12.45 & 15.71 & 来自业务模型组件加权，不是所有标的统一套 PE \\
市场情绪锚 & \multicolumn{3}{r}{53.35} & 成交额分位 +69.2\%，情绪状态 活跃共识，当前价较内在锚溢价 +512.0\% \\
综合目标价 & \multicolumn{3}{r}{53.35} & 内在/市场/券商权重：内在45\% / 市场55\% / 券商0\% \\
当前价格 & \multicolumn{3}{r}{76.22} & 2026-06-26 收盘后行情包 \\
空间与动作 & \multicolumn{3}{r}{-30.0\%；\stance{riskamber}{中性观察}} & 证据质量 B-；市场高成交已给出情绪溢价，维持中性观察并等待利润/现金流追认。 \\

\bottomrule
\end{tabularx}
\end{exhibitbox}

压力测试上，如果 2027E EPS 下修 15\%，仅 PE 组件受影响，烽火通信 的内在价值锚降至约 CNY12.09；若市场风险偏好下行导致所有估值组件倍数降 20\%，内在锚降至约 CNY9.96。两者同时发生时，内在锚约为 CNY9.67。若成交额分位回落、题材拥挤退潮或券商/市场共识下修，市场情绪锚也会同步下移，综合目标价会向内在价值锚收敛。这就是本报告要求用 Q2/Q3 交付、现金流和订单质量确认情绪溢价的原因。


\subsection{中兴通讯（000063）估值卡}
\begin{exhibitbox}[估值卡：中兴通讯]
\centering
\scriptsize
\begin{tabularx}{\exhibitboxwidth}{>{\bfseries\raggedright\arraybackslash}p{2.5cm} >{\raggedleft\arraybackslash}p{1.6cm} >{\raggedleft\arraybackslash}p{1.6cm} >{\raggedleft\arraybackslash}p{1.6cm} >{\raggedright\arraybackslash\sloppy\hspace{0pt}}X}
\toprule
\textbf{项目} & \textbf{熊市} & \textbf{基准} & \textbf{牛市} & \textbf{说明} \\
\midrule
估值方法 & \multicolumn{3}{r}{PE/PB/PS} & 55\% PE + 20\% PB + 25\% PS，用于运营商/云网络设备业务；权重 PE 55\% / PB 20\% / PS 25\% \\
EPS 分母 & 1.23 & 1.37 & 1.51 & 2026E 由 Q1 季节性折算，2027/2028E 由公司层面成长假设外推 \\
收入/股 & 32.77 & 36.70 & 40.37 & PS 组件用于校验收入规模、业务纯度和主题期权，不替代交付验证 \\
BVPS & 16.03 & 16.03 & 16.03 & PB 组件用于资产、项目和混合业务的底部校验 \\
PE 组件 & 18.41 & 27.49 & 33.75 & PE 档位 15x/20x/25x，只按权重进入综合价值 \\
PB 组件 & 22.44 & 33.67 & 40.08 & PB 档位 1.4x/2.1x/2.8x，非资产型公司为辅助或 n.a. \\
PS 组件 & 21.30 & 36.70 & 48.66 & PS 档位 0.65x/1.0x/1.35x，牛市价值按 12\% 折现 \\
内在价值锚 & 19.94 & 31.03 & 38.74 & 来自业务模型组件加权，不是所有标的统一套 PE \\
市场情绪锚 & \multicolumn{3}{r}{25.68} & 成交额分位 +57.7\%，情绪状态 活跃共识，当前价较内在锚溢价 +15.0\% \\
综合目标价 & \multicolumn{3}{r}{31.06} & 内在/市场/券商权重：内在50\% / 市场35\% / 券商15\% \\
当前价格 & \multicolumn{3}{r}{35.67} & 2026-06-26 收盘后行情包 \\
空间与动作 & \multicolumn{3}{r}{-12.9\%；\stance{riskamber}{中性}} & 证据质量 A-；动作由综合目标价相对当前价决定。 \\

\bottomrule
\end{tabularx}
\end{exhibitbox}

压力测试上，如果 2027E EPS 下修 15\%，仅 PE 组件受影响，中兴通讯 的内在价值锚降至约 CNY28.76；若市场风险偏好下行导致所有估值组件倍数降 20\%，内在锚降至约 CNY24.82。两者同时发生时，内在锚约为 CNY23.01。若成交额分位回落、题材拥挤退潮或券商/市场共识下修，市场情绪锚也会同步下移，综合目标价会向内在价值锚收敛。这就是本报告要求用 Q2/Q3 交付、现金流和订单质量确认情绪溢价的原因。


\subsection{仕佳光子（688313）估值卡}
\begin{exhibitbox}[估值卡：仕佳光子]
\centering
\scriptsize
\begin{tabularx}{\exhibitboxwidth}{>{\bfseries\raggedright\arraybackslash}p{2.5cm} >{\raggedleft\arraybackslash}p{1.6cm} >{\raggedleft\arraybackslash}p{1.6cm} >{\raggedleft\arraybackslash}p{1.6cm} >{\raggedright\arraybackslash\sloppy\hspace{0pt}}X}
\toprule
\textbf{项目} & \textbf{熊市} & \textbf{基准} & \textbf{牛市} & \textbf{说明} \\
\midrule
估值方法 & \multicolumn{3}{r}{PE/PB/PS} & 45\% PE + 15\% PB + 40\% PS，用于小规模但具战略稀缺性的光芯片、相干模块和精密光学；权重 PE 45\% / PB 15\% / PS 40\% \\
EPS 分母 & 1.12 & 1.40 & 1.65 & 2026E 由 Q1 季节性折算，2027/2028E 由公司层面成长假设外推 \\
收入/股 & 5.55 & 6.93 & 8.18 & PS 组件用于校验收入规模、业务纯度和主题期权，不替代交付验证 \\
BVPS & 3.67 & 3.67 & 3.67 & PB 组件用于资产、项目和混合业务的底部校验 \\
PE 组件 & 39.11 & 69.84 & 95.65 & PE 档位 35x/50x/65x，只按权重进入综合价值 \\
PB 组件 & 9.17 & 14.68 & 18.02 & PB 档位 2.5x/4.0x/5.5x，非资产型公司为辅助或 n.a. \\
PS 组件 & 16.64 & 34.66 & 51.12 & PS 档位 3.0x/5.0x/7.0x，牛市价值按 12\% 折现 \\
内在价值锚 & 25.63 & 47.49 & 66.20 & 来自业务模型组件加权，不是所有标的统一套 PE \\
市场情绪锚 & \multicolumn{3}{r}{125.79} & 成交额分位 +46.2\%，情绪状态 活跃共识，当前价较内在锚溢价 +278.4\% \\
综合目标价 & \multicolumn{3}{r}{125.79} & 内在/市场/券商权重：内在47\% / 市场38\% / 券商15\% \\
当前价格 & \multicolumn{3}{r}{179.70} & 2026-06-26 收盘后行情包 \\
空间与动作 & \multicolumn{3}{r}{-30.0\%；\stance{riskamber}{中性观察}} & 证据质量 B；市场高成交已给出情绪溢价，维持中性观察并等待利润/现金流追认。 \\

\bottomrule
\end{tabularx}
\end{exhibitbox}

压力测试上，如果 2027E EPS 下修 15\%，仅 PE 组件受影响，仕佳光子 的内在价值锚降至约 CNY42.78；若市场风险偏好下行导致所有估值组件倍数降 20\%，内在锚降至约 CNY37.99。两者同时发生时，内在锚约为 CNY34.22。若成交额分位回落、题材拥挤退潮或券商/市场共识下修，市场情绪锚也会同步下移，综合目标价会向内在价值锚收敛。这就是本报告要求用 Q2/Q3 交付、现金流和订单质量确认情绪溢价的原因。


\subsection{长光华芯（688048）估值卡}
\begin{exhibitbox}[估值卡：长光华芯]
\centering
\scriptsize
\begin{tabularx}{\exhibitboxwidth}{>{\bfseries\raggedright\arraybackslash}p{2.5cm} >{\raggedleft\arraybackslash}p{1.6cm} >{\raggedleft\arraybackslash}p{1.6cm} >{\raggedleft\arraybackslash}p{1.6cm} >{\raggedright\arraybackslash\sloppy\hspace{0pt}}X}
\toprule
\textbf{项目} & \textbf{熊市} & \textbf{基准} & \textbf{牛市} & \textbf{说明} \\
\midrule
估值方法 & \multicolumn{3}{r}{PE/PB/PS} & 45\% PE + 15\% PB + 40\% PS，用于小规模但具战略稀缺性的光芯片、相干模块和精密光学；权重 PE 45\% / PB 15\% / PS 40\% \\
EPS 分母 & 0.13 & 0.18 & 0.25 & 2026E 由 Q1 季节性折算，2027/2028E 由公司层面成长假设外推 \\
收入/股 & 3.68 & 5.34 & 7.21 & PS 组件用于校验收入规模、业务纯度和主题期权，不替代交付验证 \\
BVPS & 17.10 & 17.10 & 17.10 & PB 组件用于资产、项目和混合业务的底部校验 \\
PE 组件 & 5.71 & 12.89 & 19.98 & PE 档位 45x/70x/90x，只按权重进入综合价值 \\
PB 组件 & 42.74 & 68.38 & 83.95 & PB 档位 2.5x/4.0x/5.5x，非资产型公司为辅助或 n.a. \\
PS 组件 & 11.05 & 26.71 & 45.07 & PS 档位 3.0x/5.0x/7.0x，牛市价值按 12\% 折现 \\
内在价值锚 & 13.40 & 26.74 & 39.61 & 来自业务模型组件加权，不是所有标的统一套 PE \\
市场情绪锚 & \multicolumn{3}{r}{277.20} & 成交额分位 +34.6\%，情绪状态 活跃共识，当前价较内在锚溢价 +1380.9\% \\
综合目标价 & \multicolumn{3}{r}{277.20} & 内在/市场/券商权重：内在47\% / 市场38\% / 券商15\% \\
当前价格 & \multicolumn{3}{r}{396.00} & 2026-06-26 收盘后行情包 \\
空间与动作 & \multicolumn{3}{r}{-30.0\%；\stance{riskamber}{中性观察}} & 证据质量 C+；市场高成交已给出情绪溢价，维持中性观察并等待利润/现金流追认。 \\

\bottomrule
\end{tabularx}
\end{exhibitbox}

压力测试上，如果 2027E EPS 下修 15\%，仅 PE 组件受影响，长光华芯 的内在价值锚降至约 CNY25.87；若市场风险偏好下行导致所有估值组件倍数降 20\%，内在锚降至约 CNY21.39。两者同时发生时，内在锚约为 CNY20.70。若成交额分位回落、题材拥挤退潮或券商/市场共识下修，市场情绪锚也会同步下移，综合目标价会向内在价值锚收敛。这就是本报告要求用 Q2/Q3 交付、现金流和订单质量确认情绪溢价的原因。


\subsection{光库科技（300620）估值卡}
\begin{exhibitbox}[估值卡：光库科技]
\centering
\scriptsize
\begin{tabularx}{\exhibitboxwidth}{>{\bfseries\raggedright\arraybackslash}p{2.5cm} >{\raggedleft\arraybackslash}p{1.6cm} >{\raggedleft\arraybackslash}p{1.6cm} >{\raggedleft\arraybackslash}p{1.6cm} >{\raggedright\arraybackslash\sloppy\hspace{0pt}}X}
\toprule
\textbf{项目} & \textbf{熊市} & \textbf{基准} & \textbf{牛市} & \textbf{说明} \\
\midrule
估值方法 & \multicolumn{3}{r}{PE/PB/PS} & 45\% 周期正常化 PE + 30\% PB/ROE + 25\% PS，用于资产较重的光纤光缆和海缆平台；权重 PE 45\% / PB 30\% / PS 25\% \\
EPS 分母 & 0.78 & 0.98 & 1.15 & 2026E 由 Q1 季节性折算，2027/2028E 由公司层面成长假设外推 \\
收入/股 & 7.44 & 9.30 & 10.97 & PS 组件用于校验收入规模、业务纯度和主题期权，不替代交付验证 \\
BVPS & 8.85 & 8.85 & 8.85 & PB 组件用于资产、项目和混合业务的底部校验 \\
PE 组件 & 27.32 & 48.78 & 66.81 & PE 档位 35x/50x/65x，只按权重进入综合价值 \\
PB 组件 & 11.50 & 15.93 & 18.96 & PB 档位 1.3x/1.8x/2.4x，非资产型公司为辅助或 n.a. \\
PS 组件 & 4.09 & 7.44 & 10.29 & PS 档位 0.55x/0.8x/1.05x，牛市价值按 12\% 折现 \\
内在价值锚 & 16.77 & 28.59 & 38.32 & 来自业务模型组件加权，不是所有标的统一套 PE \\
市场情绪锚 & \multicolumn{3}{r}{253.70} & 成交额分位 +30.8\%，情绪状态 活跃共识，当前价较内在锚溢价 +1167.8\% \\
综合目标价 & \multicolumn{3}{r}{253.70} & 内在/市场/券商权重：内在38\% / 市场47\% / 券商15\% \\
当前价格 & \multicolumn{3}{r}{362.43} & 2026-06-26 收盘后行情包 \\
空间与动作 & \multicolumn{3}{r}{-30.0\%；\stance{riskamber}{中性观察}} & 证据质量 B；市场高成交已给出情绪溢价，维持中性观察并等待利润/现金流追认。 \\

\bottomrule
\end{tabularx}
\end{exhibitbox}

压力测试上，如果 2027E EPS 下修 15\%，仅 PE 组件受影响，光库科技 的内在价值锚降至约 CNY25.29；若市场风险偏好下行导致所有估值组件倍数降 20\%，内在锚降至约 CNY22.87。两者同时发生时，内在锚约为 CNY20.24。若成交额分位回落、题材拥挤退潮或券商/市场共识下修，市场情绪锚也会同步下移，综合目标价会向内在价值锚收敛。这就是本报告要求用 Q2/Q3 交付、现金流和订单质量确认情绪溢价的原因。


\subsection{通鼎互联（002491）估值卡}
\begin{exhibitbox}[估值卡：通鼎互联]
\centering
\scriptsize
\begin{tabularx}{\exhibitboxwidth}{>{\bfseries\raggedright\arraybackslash}p{2.5cm} >{\raggedleft\arraybackslash}p{1.6cm} >{\raggedleft\arraybackslash}p{1.6cm} >{\raggedleft\arraybackslash}p{1.6cm} >{\raggedright\arraybackslash\sloppy\hspace{0pt}}X}
\toprule
\textbf{项目} & \textbf{熊市} & \textbf{基准} & \textbf{牛市} & \textbf{说明} \\
\midrule
估值方法 & \multicolumn{3}{r}{PE/PB/PS} & 20\% 周期正常化 PE + 30\% PB/ROE + 50\% 收入期权 PS，用于光纤光缆和通信集成敞口；权重 PE 20\% / PB 30\% / PS 50\% \\
EPS 分母 & 0.18 & 0.20 & 0.22 & 2026E 由 Q1 季节性折算，2027/2028E 由公司层面成长假设外推 \\
收入/股 & 3.76 & 4.21 & 4.63 & PS 组件用于校验收入规模、业务纯度和主题期权，不替代交付验证 \\
BVPS & 2.05 & 2.05 & 2.05 & PB 组件用于资产、项目和混合业务的底部校验 \\
PE 组件 & 3.26 & 4.87 & 5.98 & PE 档位 18x/24x/30x，只按权重进入综合价值 \\
PB 组件 & 4.92 & 8.20 & 10.06 & PB 档位 2.4x/4.0x/5.5x，非资产型公司为辅助或 n.a. \\
PS 组件 & 8.27 & 21.06 & 31.03 & PS 档位 2.2x/5.0x/7.5x，牛市价值按 12\% 折现 \\
内在价值锚 & 6.27 & 13.96 & 19.73 & 来自业务模型组件加权，不是所有标的统一套 PE \\
市场情绪锚 & \multicolumn{3}{r}{26.03} & 成交额分位 +50.0\%，情绪状态 活跃共识，当前价较内在锚溢价 +158.9\% \\
综合目标价 & \multicolumn{3}{r}{26.03} & 内在/市场/券商权重：内在45\% / 市场55\% / 券商0\% \\
当前价格 & \multicolumn{3}{r}{36.15} & 2026-06-26 收盘后行情包 \\
空间与动作 & \multicolumn{3}{r}{-28.0\%；\stance{riskamber}{中性观察}} & 证据质量 B-；市场高成交已给出情绪溢价，维持中性观察并等待利润/现金流追认。 \\

\bottomrule
\end{tabularx}
\end{exhibitbox}

压力测试上，如果 2027E EPS 下修 15\%，仅 PE 组件受影响，通鼎互联 的内在价值锚降至约 CNY13.82；若市场风险偏好下行导致所有估值组件倍数降 20\%，内在锚降至约 CNY11.17。两者同时发生时，内在锚约为 CNY11.05。若成交额分位回落、题材拥挤退潮或券商/市场共识下修，市场情绪锚也会同步下移，综合目标价会向内在价值锚收敛。这就是本报告要求用 Q2/Q3 交付、现金流和订单质量确认情绪溢价的原因。


\subsection{永鼎股份（600105）估值卡}
\begin{exhibitbox}[估值卡：永鼎股份]
\centering
\scriptsize
\begin{tabularx}{\exhibitboxwidth}{>{\bfseries\raggedright\arraybackslash}p{2.5cm} >{\raggedleft\arraybackslash}p{1.6cm} >{\raggedleft\arraybackslash}p{1.6cm} >{\raggedleft\arraybackslash}p{1.6cm} >{\raggedright\arraybackslash\sloppy\hspace{0pt}}X}
\toprule
\textbf{项目} & \textbf{熊市} & \textbf{基准} & \textbf{牛市} & \textbf{说明} \\
\midrule
估值方法 & \multicolumn{3}{r}{PE/PB/PS} & 20\% 周期正常化 PE + 30\% PB/ROE + 50\% 收入期权 PS，用于光纤光缆和通信集成敞口；权重 PE 20\% / PB 30\% / PS 50\% \\
EPS 分母 & 0.82 & 0.90 & 0.97 & 2026E 由 Q1 季节性折算，2027/2028E 由公司层面成长假设外推 \\
收入/股 & 6.44 & 7.08 & 7.65 & PS 组件用于校验收入规模、业务纯度和主题期权，不替代交付验证 \\
BVPS & 2.27 & 2.27 & 2.27 & PB 组件用于资产、项目和混合业务的底部校验 \\
PE 组件 & 14.77 & 20.76 & 24.37 & PE 档位 18x/23x/28x，只按权重进入综合价值 \\
PB 组件 & 5.45 & 9.08 & 11.15 & PB 档位 2.4x/4.0x/5.5x，非资产型公司为辅助或 n.a. \\
PS 组件 & 14.16 & 35.39 & 51.20 & PS 档位 2.2x/5.0x/7.5x，牛市价值按 12\% 折现 \\
内在价值锚 & 11.67 & 24.57 & 33.82 & 来自业务模型组件加权，不是所有标的统一套 PE \\
市场情绪锚 & \multicolumn{3}{r}{46.00} & 成交额分位 +73.1\%，情绪状态 活跃共识，当前价较内在锚溢价 +167.4\% \\
综合目标价 & \multicolumn{3}{r}{46.00} & 内在/市场/券商权重：内在45\% / 市场55\% / 券商0\% \\
当前价格 & \multicolumn{3}{r}{65.71} & 2026-06-26 收盘后行情包 \\
空间与动作 & \multicolumn{3}{r}{-30.0\%；\stance{riskamber}{中性观察}} & 证据质量 B；市场高成交已给出情绪溢价，维持中性观察并等待利润/现金流追认。 \\
H1 预告(已计入) & \multicolumn{3}{r}{已入分母} & H1 2026 归母 CNY5.00--7.00 亿元、同比 +57\%--+120\%（2026-07-06）；已按 H1 年化计入 2026E EPS/内在锚，冻结对照见第 11 章 \\
\bottomrule
\end{tabularx}
\end{exhibitbox}

压力测试上，如果 2027E EPS 下修 15\%，仅 PE 组件受影响，永鼎股份 的内在价值锚降至约 CNY23.95；若市场风险偏好下行导致所有估值组件倍数降 20\%，内在锚降至约 CNY19.66。两者同时发生时，内在锚约为 CNY19.16。若成交额分位回落、题材拥挤退潮或券商/市场共识下修，市场情绪锚也会同步下移，综合目标价会向内在价值锚收敛。这就是本报告要求用 Q2/Q3 交付、现金流和订单质量确认情绪溢价的原因。


\subsection{长芯博创（300548）估值卡}
\begin{exhibitbox}[估值卡：长芯博创]
\centering
\scriptsize
\begin{tabularx}{\exhibitboxwidth}{>{\bfseries\raggedright\arraybackslash}p{2.5cm} >{\raggedleft\arraybackslash}p{1.6cm} >{\raggedleft\arraybackslash}p{1.6cm} >{\raggedleft\arraybackslash}p{1.6cm} >{\raggedright\arraybackslash\sloppy\hspace{0pt}}X}
\toprule
\textbf{项目} & \textbf{熊市} & \textbf{基准} & \textbf{牛市} & \textbf{说明} \\
\midrule
估值方法 & \multicolumn{3}{r}{PE/PB/PS} & 70\% PE + 10\% PB + 20\% PS，用于有正 EPS 但 AI 纯度不完美的模块/器件平台；权重 PE 70\% / PB 10\% / PS 20\% \\
EPS 分母 & 1.88 & 2.34 & 2.77 & 2026E 由 Q1 季节性折算，2027/2028E 由公司层面成长假设外推 \\
收入/股 & 9.67 & 12.08 & 14.26 & PS 组件用于校验收入规模、业务纯度和主题期权，不替代交付验证 \\
BVPS & 7.17 & 7.17 & 7.17 & PB 组件用于资产、项目和混合业务的底部校验 \\
PE 组件 & 56.25 & 98.44 & 135.81 & PE 档位 30x/42x/55x，只按权重进入综合价值 \\
PB 组件 & 17.93 & 28.68 & 32.01 & PB 档位 2.5x/4.0x/5.0x，非资产型公司为辅助或 n.a. \\
PS 组件 & 19.33 & 42.29 & 63.66 & PS 档位 2.0x/3.5x/5.0x，牛市价值按 12\% 折现 \\
内在价值锚 & 45.03 & 80.23 & 111.00 & 来自业务模型组件加权，不是所有标的统一套 PE \\
市场情绪锚 & \multicolumn{3}{r}{189.11} & 成交额分位 +38.5\%，情绪状态 活跃共识，当前价较内在锚溢价 +227.4\% \\
综合目标价 & \multicolumn{3}{r}{189.11} & 内在/市场/券商权重：内在55\% / 市场45\% / 券商0\% \\
当前价格 & \multicolumn{3}{r}{262.65} & 2026-06-26 收盘后行情包 \\
空间与动作 & \multicolumn{3}{r}{-28.0\%；\stance{riskamber}{中性观察}} & 证据质量 B；市场高成交已给出情绪溢价，维持中性观察并等待利润/现金流追认。 \\

\bottomrule
\end{tabularx}
\end{exhibitbox}

压力测试上，如果 2027E EPS 下修 15\%，仅 PE 组件受影响，长芯博创 的内在价值锚降至约 CNY69.90；若市场风险偏好下行导致所有估值组件倍数降 20\%，内在锚降至约 CNY64.19。两者同时发生时，内在锚约为 CNY55.92。若成交额分位回落、题材拥挤退潮或券商/市场共识下修，市场情绪锚也会同步下移，综合目标价会向内在价值锚收敛。这就是本报告要求用 Q2/Q3 交付、现金流和订单质量确认情绪溢价的原因。


\subsection{德科立（688205）估值卡}
\begin{exhibitbox}[估值卡：德科立]
\centering
\scriptsize
\begin{tabularx}{\exhibitboxwidth}{>{\bfseries\raggedright\arraybackslash}p{2.5cm} >{\raggedleft\arraybackslash}p{1.6cm} >{\raggedleft\arraybackslash}p{1.6cm} >{\raggedleft\arraybackslash}p{1.6cm} >{\raggedright\arraybackslash\sloppy\hspace{0pt}}X}
\toprule
\textbf{项目} & \textbf{熊市} & \textbf{基准} & \textbf{牛市} & \textbf{说明} \\
\midrule
估值方法 & \multicolumn{3}{r}{PE/PB/PS} & 45\% PE + 15\% PB + 40\% PS，用于小规模但具战略稀缺性的光芯片、相干模块和精密光学；权重 PE 45\% / PB 15\% / PS 40\% \\
EPS 分母 & 0.55 & 0.68 & 0.80 & 2026E 由 Q1 季节性折算，2027/2028E 由公司层面成长假设外推 \\
收入/股 & 7.07 & 8.76 & 10.34 & PS 组件用于校验收入规模、业务纯度和主题期权，不替代交付验证 \\
BVPS & 14.79 & 14.79 & 14.79 & PB 组件用于资产、项目和混合业务的底部校验 \\
PE 组件 & 16.36 & 30.44 & 41.33 & PE 档位 30x/45x/58x，只按权重进入综合价值 \\
PB 组件 & 36.97 & 59.16 & 72.63 & PB 档位 2.5x/4.0x/5.5x，非资产型公司为辅助或 n.a. \\
PS 组件 & 21.20 & 43.82 & 64.63 & PS 档位 3.0x/5.0x/7.0x，牛市价值按 12\% 折现 \\
内在价值锚 & 21.39 & 40.10 & 55.34 & 来自业务模型组件加权，不是所有标的统一套 PE \\
市场情绪锚 & \multicolumn{3}{r}{134.73} & 成交额分位 +15.4\%，情绪状态 中性交易，当前价较内在锚溢价 +425.0\% \\
综合目标价 & \multicolumn{3}{r}{73.50} & 内在/市场/券商权重：内在65\% / 市场35\% / 券商0\% \\
当前价格 & \multicolumn{3}{r}{210.52} & 2026-06-26 收盘后行情包 \\
空间与动作 & \multicolumn{3}{r}{-65.1\%；\stance{riskred}{减持}} & 证据质量 B；动作由综合目标价相对当前价决定。 \\

\bottomrule
\end{tabularx}
\end{exhibitbox}

压力测试上，如果 2027E EPS 下修 15\%，仅 PE 组件受影响，德科立 的内在价值锚降至约 CNY38.04；若市场风险偏好下行导致所有估值组件倍数降 20\%，内在锚降至约 CNY32.08。两者同时发生时，内在锚约为 CNY30.43。若成交额分位回落、题材拥挤退潮或券商/市场共识下修，市场情绪锚也会同步下移，综合目标价会向内在价值锚收敛。这就是本报告要求用 Q2/Q3 交付、现金流和订单质量确认情绪溢价的原因。


\subsection{腾景科技（688195）估值卡}
\begin{exhibitbox}[估值卡：腾景科技]
\centering
\scriptsize
\begin{tabularx}{\exhibitboxwidth}{>{\bfseries\raggedright\arraybackslash}p{2.5cm} >{\raggedleft\arraybackslash}p{1.6cm} >{\raggedleft\arraybackslash}p{1.6cm} >{\raggedleft\arraybackslash}p{1.6cm} >{\raggedright\arraybackslash\sloppy\hspace{0pt}}X}
\toprule
\textbf{项目} & \textbf{熊市} & \textbf{基准} & \textbf{牛市} & \textbf{说明} \\
\midrule
估值方法 & \multicolumn{3}{r}{PE/PB/PS} & 45\% PE + 15\% PB + 40\% PS，用于小规模但具战略稀缺性的光芯片、相干模块和精密光学；权重 PE 45\% / PB 15\% / PS 40\% \\
EPS 分母 & 0.48 & 0.57 & 0.66 & 2026E 由 Q1 季节性折算，2027/2028E 由公司层面成长假设外推 \\
收入/股 & 5.67 & 6.81 & 7.83 & PS 组件用于校验收入规模、业务纯度和主题期权，不替代交付验证 \\
BVPS & 7.87 & 7.87 & 7.87 & PB 组件用于资产、项目和混合业务的底部校验 \\
PE 组件 & 14.35 & 24.10 & 32.41 & PE 档位 30x/42x/55x，只按权重进入综合价值 \\
PB 组件 & 19.67 & 31.48 & 38.64 & PB 档位 2.5x/4.0x/5.5x，非资产型公司为辅助或 n.a. \\
PS 组件 & 17.02 & 34.03 & 48.92 & PS 档位 3.0x/5.0x/7.0x，牛市价值按 12\% 折现 \\
内在价值锚 & 16.21 & 29.18 & 39.95 & 来自业务模型组件加权，不是所有标的统一套 PE \\
市场情绪锚 & \multicolumn{3}{r}{112.56} & 成交额分位 +19.2\%，情绪状态 中性交易，当前价较内在锚溢价 +542.9\% \\
综合目标价 & \multicolumn{3}{r}{83.18} & 内在/市场/券商权重：内在55\% / 市场30\% / 券商15\% \\
当前价格 & \multicolumn{3}{r}{187.60} & 2026-06-26 收盘后行情包 \\
空间与动作 & \multicolumn{3}{r}{-55.7\%；\stance{riskred}{减持}} & 证据质量 B；动作由综合目标价相对当前价决定。 \\

\bottomrule
\end{tabularx}
\end{exhibitbox}

压力测试上，如果 2027E EPS 下修 15\%，仅 PE 组件受影响，腾景科技 的内在价值锚降至约 CNY27.55；若市场风险偏好下行导致所有估值组件倍数降 20\%，内在锚降至约 CNY23.34。两者同时发生时，内在锚约为 CNY22.04。若成交额分位回落、题材拥挤退潮或券商/市场共识下修，市场情绪锚也会同步下移，综合目标价会向内在价值锚收敛。这就是本报告要求用 Q2/Q3 交付、现金流和订单质量确认情绪溢价的原因。


\subsection{联特科技（301205）估值卡}
\begin{exhibitbox}[估值卡：联特科技]
\centering
\scriptsize
\begin{tabularx}{\exhibitboxwidth}{>{\bfseries\raggedright\arraybackslash}p{2.5cm} >{\raggedleft\arraybackslash}p{1.6cm} >{\raggedleft\arraybackslash}p{1.6cm} >{\raggedleft\arraybackslash}p{1.6cm} >{\raggedright\arraybackslash\sloppy\hspace{0pt}}X}
\toprule
\textbf{项目} & \textbf{熊市} & \textbf{基准} & \textbf{牛市} & \textbf{说明} \\
\midrule
估值方法 & \multicolumn{3}{r}{PE/PB/PS} & 45\% PE + 15\% PB + 40\% PS，用于小规模但具战略稀缺性的光芯片、相干模块和精密光学；权重 PE 45\% / PB 15\% / PS 40\% \\
EPS 分母 & 0.12 & 0.16 & 0.20 & 2026E 由 Q1 季节性折算，2027/2028E 由公司层面成长假设外推 \\
收入/股 & 8.23 & 11.11 & 13.89 & PS 组件用于校验收入规模、业务纯度和主题期权，不替代交付验证 \\
BVPS & 12.55 & 12.55 & 12.55 & PB 组件用于资产、项目和混合业务的底部校验 \\
PE 组件 & 4.11 & 8.72 & 12.39 & PE 档位 35x/55x/70x，只按权重进入综合价值 \\
PB 组件 & 31.37 & 50.19 & 61.62 & PB 档位 2.5x/4.0x/5.5x，非资产型公司为辅助或 n.a. \\
PS 组件 & 24.69 & 55.56 & 86.81 & PS 档位 3.0x/5.0x/7.0x，牛市价值按 12\% 折现 \\
内在价值锚 & 16.43 & 33.68 & 49.54 & 来自业务模型组件加权，不是所有标的统一套 PE \\
市场情绪锚 & \multicolumn{3}{r}{189.52} & 成交额分位 +23.1\%，情绪状态 中性交易，当前价较内在锚溢价 +837.9\% \\
综合目标价 & \multicolumn{3}{r}{116.17} & 内在/市场/券商权重：内在55\% / 市场30\% / 券商15\% \\
当前价格 & \multicolumn{3}{r}{315.86} & 2026-06-26 收盘后行情包 \\
空间与动作 & \multicolumn{3}{r}{-63.2\%；\stance{riskred}{减持}} & 证据质量 C+；动作由综合目标价相对当前价决定。 \\

\bottomrule
\end{tabularx}
\end{exhibitbox}

压力测试上，如果 2027E EPS 下修 15\%，仅 PE 组件受影响，联特科技 的内在价值锚降至约 CNY33.09；若市场风险偏好下行导致所有估值组件倍数降 20\%，内在锚降至约 CNY26.94。两者同时发生时，内在锚约为 CNY26.47。若成交额分位回落、题材拥挤退潮或券商/市场共识下修，市场情绪锚也会同步下移，综合目标价会向内在价值锚收敛。这就是本报告要求用 Q2/Q3 交付、现金流和订单质量确认情绪溢价的原因。


\subsection{兆龙互连（300913）估值卡}
\begin{exhibitbox}[估值卡：兆龙互连]
\centering
\scriptsize
\begin{tabularx}{\exhibitboxwidth}{>{\bfseries\raggedright\arraybackslash}p{2.5cm} >{\raggedleft\arraybackslash}p{1.6cm} >{\raggedleft\arraybackslash}p{1.6cm} >{\raggedleft\arraybackslash}p{1.6cm} >{\raggedright\arraybackslash\sloppy\hspace{0pt}}X}
\toprule
\textbf{项目} & \textbf{熊市} & \textbf{基准} & \textbf{牛市} & \textbf{说明} \\
\midrule
估值方法 & \multicolumn{3}{r}{PE/PB/PS} & 40\% PE + 25\% PB + 35\% PS，用于连接器、线缆和高速互连混合业务；权重 PE 40\% / PB 25\% / PS 35\% \\
EPS 分母 & 0.59 & 0.70 & 0.79 & 2026E 由 Q1 季节性折算，2027/2028E 由公司层面成长假设外推 \\
收入/股 & 6.18 & 7.29 & 8.31 & PS 组件用于校验收入规模、业务纯度和主题期权，不替代交付验证 \\
BVPS & 7.73 & 7.73 & 7.73 & PB 组件用于资产、项目和混合业务的底部校验 \\
PE 组件 & 14.75 & 24.37 & 31.89 & PE 档位 25x/35x/45x，只按权重进入综合价值 \\
PB 组件 & 15.46 & 27.05 & 33.12 & PB 档位 2.0x/3.5x/4.8x，非资产型公司为辅助或 n.a. \\
PS 组件 & 5.56 & 13.13 & 20.78 & PS 档位 0.9x/1.8x/2.8x，牛市价值按 12\% 折现 \\
内在价值锚 & 11.71 & 21.10 & 28.31 & 来自业务模型组件加权，不是所有标的统一套 PE \\
市场情绪锚 & \multicolumn{3}{r}{24.36} & 成交额分位 +7.7\%，情绪状态 弱共识，当前价较内在锚溢价 +140.5\% \\
综合目标价 & \multicolumn{3}{r}{22.44} & 内在/市场/券商权重：内在59\% / 市场41\% / 券商0\% \\
当前价格 & \multicolumn{3}{r}{50.75} & 2026-06-26 收盘后行情包 \\
空间与动作 & \multicolumn{3}{r}{-55.8\%；\stance{riskred}{减持}} & 证据质量 B；动作由综合目标价相对当前价决定。 \\

\bottomrule
\end{tabularx}
\end{exhibitbox}

压力测试上，如果 2027E EPS 下修 15\%，仅 PE 组件受影响，兆龙互连 的内在价值锚降至约 CNY19.64；若市场风险偏好下行导致所有估值组件倍数降 20\%，内在锚降至约 CNY16.88。两者同时发生时，内在锚约为 CNY15.71。若成交额分位回落、题材拥挤退潮或券商/市场共识下修，市场情绪锚也会同步下移，综合目标价会向内在价值锚收敛。这就是本报告要求用 Q2/Q3 交付、现金流和订单质量确认情绪溢价的原因。


\subsection{意华股份（002897）估值卡}
\begin{exhibitbox}[估值卡：意华股份]
\centering
\scriptsize
\begin{tabularx}{\exhibitboxwidth}{>{\bfseries\raggedright\arraybackslash}p{2.5cm} >{\raggedleft\arraybackslash}p{1.6cm} >{\raggedleft\arraybackslash}p{1.6cm} >{\raggedleft\arraybackslash}p{1.6cm} >{\raggedright\arraybackslash\sloppy\hspace{0pt}}X}
\toprule
\textbf{项目} & \textbf{熊市} & \textbf{基准} & \textbf{牛市} & \textbf{说明} \\
\midrule
估值方法 & \multicolumn{3}{r}{PE/PB/PS} & 40\% PE + 25\% PB + 35\% PS，用于连接器、线缆和高速互连混合业务；权重 PE 40\% / PB 25\% / PS 35\% \\
EPS 分母 & 0.95 & 1.10 & 1.23 & 2026E 由 Q1 季节性折算，2027/2028E 由公司层面成长假设外推 \\
收入/股 & 30.64 & 35.23 & 39.46 & PS 组件用于校验收入规模、业务纯度和主题期权，不替代交付验证 \\
BVPS & 14.66 & 14.66 & 14.66 & PB 组件用于资产、项目和混合业务的底部校验 \\
PE 组件 & 21.00 & 32.93 & 41.71 & PE 档位 22x/30x/38x，只按权重进入综合价值 \\
PB 组件 & 29.32 & 51.32 & 62.84 & PB 档位 2.0x/3.5x/4.8x，非资产型公司为辅助或 n.a. \\
PS 组件 & 27.57 & 63.42 & 98.65 & PS 档位 0.9x/1.8x/2.8x，牛市价值按 12\% 折现 \\
内在价值锚 & 25.38 & 48.20 & 66.92 & 来自业务模型组件加权，不是所有标的统一套 PE \\
市场情绪锚 & \multicolumn{3}{r}{52.39} & 成交额分位 +11.5\%，情绪状态 中性交易，当前价较内在锚溢价 +75.3\% \\
综合目标价 & \multicolumn{3}{r}{49.92} & 内在/市场/券商权重：内在59\% / 市场41\% / 券商0\% \\
当前价格 & \multicolumn{3}{r}{84.50} & 2026-06-26 收盘后行情包 \\
空间与动作 & \multicolumn{3}{r}{-40.9\%；\stance{riskred}{减持}} & 证据质量 B-；动作由综合目标价相对当前价决定。 \\

\bottomrule
\end{tabularx}
\end{exhibitbox}

压力测试上，如果 2027E EPS 下修 15\%，仅 PE 组件受影响，意华股份 的内在价值锚降至约 CNY46.22；若市场风险偏好下行导致所有估值组件倍数降 20\%，内在锚降至约 CNY38.56。两者同时发生时，内在锚约为 CNY36.98。若成交额分位回落、题材拥挤退潮或券商/市场共识下修，市场情绪锚也会同步下移，综合目标价会向内在价值锚收敛。这就是本报告要求用 Q2/Q3 交付、现金流和订单质量确认情绪溢价的原因。


\subsection{神宇股份（300563）估值卡}
\begin{exhibitbox}[估值卡：神宇股份]
\centering
\scriptsize
\begin{tabularx}{\exhibitboxwidth}{>{\bfseries\raggedright\arraybackslash}p{2.5cm} >{\raggedleft\arraybackslash}p{1.6cm} >{\raggedleft\arraybackslash}p{1.6cm} >{\raggedleft\arraybackslash}p{1.6cm} >{\raggedright\arraybackslash\sloppy\hspace{0pt}}X}
\toprule
\textbf{项目} & \textbf{熊市} & \textbf{基准} & \textbf{牛市} & \textbf{说明} \\
\midrule
估值方法 & \multicolumn{3}{r}{PE/PB/PS} & 40\% PE + 25\% PB + 35\% PS，用于连接器、线缆和高速互连混合业务；权重 PE 40\% / PB 25\% / PS 35\% \\
EPS 分母 & 0.32 & 0.37 & 0.41 & 2026E 由 Q1 季节性折算，2027/2028E 由公司层面成长假设外推 \\
收入/股 & 4.80 & 5.52 & 6.18 & PS 组件用于校验收入规模、业务纯度和主题期权，不替代交付验证 \\
BVPS & 6.33 & 6.33 & 6.33 & PB 组件用于资产、项目和混合业务的底部校验 \\
PE 组件 & 7.00 & 11.71 & 14.64 & PE 档位 22x/32x/40x，只按权重进入综合价值 \\
PB 组件 & 12.66 & 22.15 & 27.12 & PB 档位 2.0x/3.5x/4.8x，非资产型公司为辅助或 n.a. \\
PS 组件 & 4.32 & 9.94 & 15.46 & PS 档位 0.9x/1.8x/2.8x，牛市价值按 12\% 折现 \\
内在价值锚 & 7.48 & 13.70 & 18.04 & 来自业务模型组件加权，不是所有标的统一套 PE \\
市场情绪锚 & \multicolumn{3}{r}{11.87} & 成交额分位 +3.8\%，情绪状态 弱共识，当前价较内在锚溢价 +80.5\% \\
综合目标价 & \multicolumn{3}{r}{12.95} & 内在/市场/券商权重：内在59\% / 市场41\% / 券商0\% \\
当前价格 & \multicolumn{3}{r}{24.73} & 2026-06-26 收盘后行情包 \\
空间与动作 & \multicolumn{3}{r}{-47.7\%；\stance{riskred}{减持}} & 证据质量 C+；动作由综合目标价相对当前价决定。 \\

\bottomrule
\end{tabularx}
\end{exhibitbox}

压力测试上，如果 2027E EPS 下修 15\%，仅 PE 组件受影响，神宇股份 的内在价值锚降至约 CNY13.00；若市场风险偏好下行导致所有估值组件倍数降 20\%，内在锚降至约 CNY10.96。两者同时发生时，内在锚约为 CNY10.40。若成交额分位回落、题材拥挤退潮或券商/市场共识下修，市场情绪锚也会同步下移，综合目标价会向内在价值锚收敛。这就是本报告要求用 Q2/Q3 交付、现金流和订单质量确认情绪溢价的原因。


\subsection{杭电股份（603618）估值卡}
\begin{exhibitbox}[估值卡：杭电股份]
\centering
\scriptsize
\begin{tabularx}{\exhibitboxwidth}{>{\bfseries\raggedright\arraybackslash}p{2.5cm} >{\raggedleft\arraybackslash}p{1.6cm} >{\raggedleft\arraybackslash}p{1.6cm} >{\raggedleft\arraybackslash}p{1.6cm} >{\raggedright\arraybackslash\sloppy\hspace{0pt}}X}
\toprule
\textbf{项目} & \textbf{熊市} & \textbf{基准} & \textbf{牛市} & \textbf{说明} \\
\midrule
估值方法 & \multicolumn{3}{r}{PE/PB/PS} & 45\% 周期正常化 PE + 30\% PB/ROE + 25\% PS，用于资产较重的光纤光缆和海缆平台；权重 PE 45\% / PB 30\% / PS 25\% \\
EPS 分母 & 1.13 & 1.30 & 1.43 & 2026E 由 Q1 季节性折算，2027/2028E 由公司层面成长假设外推 \\
收入/股 & 30.73 & 35.34 & 38.88 & PS 组件用于校验收入规模、业务纯度和主题期权，不替代交付验证 \\
BVPS & 3.91 & 3.91 & 3.91 & PB 组件用于资产、项目和混合业务的底部校验 \\
PE 组件 & 18.05 & 28.54 & 35.68 & PE 档位 16x/22x/28x，只按权重进入综合价值 \\
PB 组件 & 5.09 & 7.04 & 8.39 & PB 档位 1.3x/1.8x/2.4x，非资产型公司为辅助或 n.a. \\
PS 组件 & 16.90 & 28.28 & 36.45 & PS 档位 0.55x/0.8x/1.05x，牛市价值按 12\% 折现 \\
内在价值锚 & 13.87 & 22.03 & 27.68 & 来自业务模型组件加权，不是所有标的统一套 PE \\
市场情绪锚 & \multicolumn{3}{r}{31.20} & 成交额分位 +26.9\%，情绪状态 中性交易，当前价较内在锚溢价 +136.1\% \\
综合目标价 & \multicolumn{3}{r}{25.80} & 内在/市场/券商权重：内在59\% / 市场41\% / 券商0\% \\
当前价格 & \multicolumn{3}{r}{52.00} & 2026-06-26 收盘后行情包 \\
空间与动作 & \multicolumn{3}{r}{-50.4\%；\stance{riskred}{减持}} & 证据质量 C+；动作由综合目标价相对当前价决定。 \\
H1 预告(已计入) & \multicolumn{3}{r}{已入分母} & H1 2026 归母 CNY3.60--4.00 亿元、同比 +852\%--+958\%（2026-07-04）；已按 H1 年化计入 2026E EPS/内在锚，冻结对照见第 11 章 \\
\bottomrule
\end{tabularx}
\end{exhibitbox}

压力测试上，如果 2027E EPS 下修 15\%，仅 PE 组件受影响，杭电股份 的内在价值锚降至约 CNY20.10；若市场风险偏好下行导致所有估值组件倍数降 20\%，内在锚降至约 CNY17.62。两者同时发生时，内在锚约为 CNY16.08。若成交额分位回落、题材拥挤退潮或券商/市场共识下修，市场情绪锚也会同步下移，综合目标价会向内在价值锚收敛。这就是本报告要求用 Q2/Q3 交付、现金流和订单质量确认情绪溢价的原因。
